\documentclass{beamer}
\usepackage{bm}
\usepackage{tabularx}
\setbeamertemplate{caption}[numbered]
\usepackage{xcolor}
\usepackage{multirow}
\usepackage{movie15}
\usepackage{Sweave}
%\usepackage{nameref}% http://ctan.org/pkg/nameref
%\newcommand{\currentchapter}{}
%\let\oldchapter\chapter
%\renewcommand{\chapter}[1]{\oldchapter{#1}\renewcommand{\currentchapter}{#1}}
\renewcommand{\thesection}{3.\arabic{section}}
\newcommand{\currentsection}{}
\let\oldsection\section
\renewcommand{\section}[1]{\oldsection{#1}\renewcommand{\currentsection}{#1}}

\newcommand{\currentsubsection}{}
\let\oldsubsection\subsection
\renewcommand{\subsection}[1]{\oldsubsection{#1}\renewcommand{\currentsubsection}{#1}}


\newcommand{\boldbeta}{{\boldsymbol \beta}}
\newcommand{\boldtheta}{{\boldsymbol \theta}}
\newcommand{\boldgamma}{{\boldsymbol \gamma}}
\newcommand{\boldmu}{{\boldsymbol \mu}}
\newcommand{\boldeta}{{\boldsymbol \eta}}
\newcommand{\boldxi}{{\boldsymbol \xi}}
\newcommand{\boldpi}{{\boldsymbol \pi}}
\newcommand{\ringLambda}{\mathring{\Lambda}}
\newcommand{\currentsubsubsection}{}
\let\oldsubsubsection\subsubsection
\renewcommand{\subsubsection}[1]{\oldsubsubsection{#1}\renewcommand{\currentsubsubsection}{#1}}

\newenvironment{wideitemize}%
  {\begin{itemize}%
%    \setlength{\itemsep}{0pt}%
    \setlength{\parskip}{1em}}%
  {\end{itemize}}


%\usetheme{CambridgeUS}
\usetheme{Madrid}

%\title[Section~1.\insertsectionnumber \currentsection]{Chapter 1. Motivating Examples}
\title[\currentsection]{Chapter 3. Models for Multi-categorical Responses: Multivariate Extensions of GLM }
\author[Multivariate GLM]{MAST90084 Statistical Modelling Slides}
\vskip 1cm
\institute[Ch3]{Dennis Leung \\ \vskip 0.3cm SCHOOL OF MATHEMATICS AND STATISTICS\\
\vskip 0.3cm
   THE UNIVERSITY OF MELBOURNE}
\vskip 1cm
\date[Dennis Leung \hspace{8pt}]{}


\begin{document}

\frame{\titlepage}
%
%\begin{frame} {Something about the first assignment}
%\end{frame}


\begin{frame}[fragile] \frametitle{Outline}
{\scriptsize
\tableofcontents}
\end{frame}

%
\section{\S3.3 Models for ordinal responses} \label{sec:3.3}
%%\begin{frame}{\S3.3 Models for ordinal responses}
%%Response variable having $k>2$ categories may be of ordinal nature.
%%
%%\begin{itemize}
%%
%%\item \textbf{Example 3.2} Breathing test results
%%
%%\includegraphics[width=.5\textwidth]{BreathingTest}
%%
%%
%%\item  \textbf{Grouped continuous} (Anderson 1984) --- merely a categorized continuous variable since the responses ``Normal", ``Borderline", etc.. can arise as intervals of a physical scale like volume of breath. 
%%\end{itemize}
%%
%%%\begin{itemize}
%%%\item
%%%Response variable having $k>2$ categories may be of ordinal nature.
%%%\begin{itemize}
%%%\item
%%%\textbf{Example 3.2} Breathing test results
%%%
%%%\includegraphics[width=.5\textwidth]{Breathingtest.png}
%%%
%%%\item
%%%\textbf{Grouped continuous} --- merely a categorized continuous variable
%%%
%%%\item
%%%\textbf{Example 3.3} Job expectations
%%%\end{itemize}
%%%\item
%%%May  come from quite different mechanisms. Two distinguished types:
%%%\textbf{grouped continuous} and \textbf{assessed ordered} categorical variables.
%%%
%%%
%%%\item
%%%\textbf{Assessed ordered} --- obtained after an assessor's judgment.
%%%\end{itemize}
%%%\end{itemize}
%%\end{frame}
%%
%%\begin{frame}
%%\begin{itemize}
%%\item \textbf{Example 3.3} Reponse: Job expectations
%%\begin{itemize}
%%
%%\item Category 1 = don't expect adequate employment, 
%%\item Category 2 = not sure 
%%\item Category 3 =  immediately after getting the degree
%%
%%\end{itemize}
%%\includegraphics[width=.5\textwidth]{JobExpectation}
%%
%%
%%\item \textbf{Assessed ordered} --- obtained after an assessor's judgment.
%%\end{itemize}
%%
%%\end{frame}
%%
%%%\begin{frame}
%%%
%%%
%%%\end{frame}
%%
%%\subsection{\S3.3.1 Cumulative (or threshold) models} \label{sec:3.3.1}
%%\begin{frame}{\S3.3.1 Cumulative models: the threshold approach}
%%\begin{itemize}
%%\item
%%Assume: $Y$, with $k$ categories, is a categorized version  of a latent continuous variable $U$
%%% for the response variable $Y$ having
%%%$k$ categories.
%%\item
%%$Y$ determined by $U$ as \vspace*{-3pt}
%%$$Y=r \quad\Leftrightarrow\quad \theta_{r-1}<U\leq \theta_r; \quad r=1,\cdots,k \vspace*{-3pt}$$ 
%%where the {\bf thresholds} $-\infty=\theta_0<\theta_1<\cdots <\theta_k=+\infty$ unknown.
%%
%%%That means $Y$ is a coarser (categorized) version of $U$ determined by the thresholds 
%%%$\theta_1, \cdots, \theta_{k-1}$.
%%\item
%%$U$ depends on a $p\times 1$ vector explanatory variable $\textbf{x}$ as \vspace*{-5pt}
%%$$U=-\textbf{x}^T\mbox{\boldmath $\gamma$}+\varepsilon$$
%%where $\mbox{\boldmath $\gamma$}=(\gamma_1,\cdots,\gamma_p)^T$ is an unknown vector
%%parameter and $\varepsilon$ is random with cdf $F$.
%%
%%\end{itemize}
%%\end{frame}
%%
%%\begin{frame}
%%It follows that \vspace*{-5pt}
%%$$P(Y\leq r |\textbf{x})=P(U\leq \theta_r)=F(\theta_r+\textbf{x}^T\mbox{\boldmath $\gamma$}). \vspace*{-3pt}$$
%%This is called a \textbf{cumulative model} or \textbf{threshold model} with cdf $F$.
%%
%%\end{frame}
%%
%%\begin{frame}{\S3.3.1 Cumulative logistic (or proportional odds) model (1)}
%%\begin{itemize}
%%
%%
%%
%%\item
%%If choose $\displaystyle F(x)=\left[1+e^{-x}\right]^{-1}$ which is the cdf of a logistic distribution, then
%%\begin{equation}
%%P(Y\leq r|\textbf{x})=\frac{e^{\theta_r+\textbf{x}^T\mbox{\boldmath $\gamma$}}}{1 
%%+e^{\theta_r+\textbf{x}^T\mbox{\boldmath $\gamma$}}}, \quad r=1,\cdots, q=k-1.
%%\label{eq-3.3.4}
%%\end{equation}
%%\item
%%It is equivalent to
%%\begin{eqnarray}
%%\text{logit} ( P(Y\leq r|\textbf{x})) &=& \log\frac{P(Y\leq r|\textbf{x})}{P(Y> r|\textbf{x})} = \theta_r+\textbf{x}^T\mbox{\boldmath $\gamma$}
%%\quad\quad \mbox{or} \label{eq-3.3.5} \\
%%\frac{P(Y\leq r|\textbf{x})}{P(Y> r|\textbf{x})} &=& \exp\{\theta_r+\textbf{x}^T\mbox{\boldmath $\gamma$}\}
%%\label{eq-3.3.6}
%%\end{eqnarray}
%%The model given by (\ref{eq-3.3.4}), (\ref{eq-3.3.5}) or (\ref{eq-3.3.6}) is called the \textbf{cumulative
%%logistic model}.
%%\end{itemize}
%%\end{frame}
%%
%%\begin{frame}{\S3.3.1 Cumulative logistic (or proportional odds) model (2)}
%%\begin{itemize}
%%
%%\item
%%Suppose there are 2 populations (groups) identified by $\textbf{x}_1$ and $\textbf{x}_2$ respectively. Then
%%$$\frac{P(Y\leq r|\textbf{x}_1)}{P(Y> r|\textbf{x}_1)} = e^{\theta_r+\textbf{x}_1^T\mbox{\boldmath $\gamma$}}
%%\quad\mbox{and}\quad \frac{P(Y\leq r|\textbf{x}_2)}{P(Y> r|\textbf{x}_2)} = e^{\theta_r+
%%\textbf{x}_2^T\mbox{\boldmath $\gamma$}}.$$
%%It follows that
%%$$\mbox{\bf cumulative odds ratio}=\frac{P(Y\leq r|\textbf{x}_1)/P(Y> r|\textbf{x}_1)}{P(Y\leq r|\textbf{x}_2)/
%%P(Y> r|\textbf{x}_2)}=e^{(\textbf{x}_1-\textbf{x}_2)^T\mbox{\boldmath $\gamma$}}$$
%%which is the same for all $r=1,\cdots,q$.
%%\end{itemize}
%%\end{frame}
%%
%%\begin{frame}
%%\begin{itemize}
%%\item
%%The same proportionality constant $\boldsymbol \gamma$ applies to all the logits. The log cumulative odds ratio is {\bf proportional} to the distance ${\bf x}_1 - {\bf x}_2$. Hence the model's also called a \textbf{proportional odds model} (McCullagh 1980). 
%%
%%\item
%%The latent variable isn't absolutely necessary for the interpretation based on (1)-(3) alone. 
%%\end{itemize}
%%\end{frame}
%%
%%\begin{frame}{\S3.3.1 Grouped Cox (or proportional hazards) model}
%%\begin{itemize}
%%\item
%%If choose $F$ as the cdf of the extreme \emph{minimal}-value distribution, i.e. 
%%$\displaystyle F(x)=1-\exp\{-e^x\}$, then
%%\begin{eqnarray}
%%P(Y\leq r|\textbf{x}) &=& 1-\exp\{-e^{\theta_r+\textbf{x}^T\mbox{\boldmath $\gamma$}}\}
%%\label{eq-3.3.8} \\   \Longleftrightarrow \quad
%%\underbrace{\log\left[-\log P(Y>r|\textbf{x})\right]}_{\mbox{complementary log-log link}}
%%&=& \theta_r+\textbf{x}^T\mbox{\boldmath $\gamma$}, \quad r=1,\cdots,q
%%\nonumber
%%\end{eqnarray}
%%Model (\ref{eq-3.3.8}) is referred to as the \textbf{grouped Cox model} or \textbf{proportional 
%%hazards model} due to the following interpretation:
%%
%%\end{itemize}
%%\end{frame}
%%
%%\begin{frame}
%%\begin{itemize}
%%\item
%%For 2 populations (groups) identified by $\textbf{x}_1$ and $\textbf{x}_2$ respectively
%%$$\log P(Y> r|\textbf{x}_1)=\log P(Y> r|\textbf{x}_2)\cdot
%%e^{(\textbf{x}_1-\textbf{x}_2)^T\mbox{\boldmath $\gamma$}}, \quad r=1,\cdots, q$$
%%Differentiating w.r.t. $r$ on both sides results in
%%$$\lambda(r|\textbf{x}_1)=\lambda(r|\textbf{x}_2)\cdot
%%e^{(\textbf{x}_1-\textbf{x}_2)^T\mbox{\boldmath $\gamma$}}, \quad \Leftrightarrow\quad
%%\frac{\lambda(r|\textbf{x}_1)}{\lambda(r|\textbf{x}_2)}=
%%e^{(\textbf{x}_1-\textbf{x}_2)^T\mbox{\boldmath $\gamma$}},$$
%%which is the same for all $r=1,\cdots,q$. Here {\footnotesize $\lambda(r|\cdot)=-\frac{d}{dr}\log P(Y>r|\cdot)$} is
%%the \textbf{hazard function}.
%%
%%\item (Will get to proportional hazard model if time allows, ch 9 of F\& T)
%%\end{itemize}
%%
%%\end{frame}
%%
%%\begin{frame}{\large \S3.3.1 Cumulative model based on extreme maximal-value CDF}
%%\begin{itemize}
%%\item
%%If choose $F$ as the cdf of the extreme \emph{maximal}-value distribution, i.e. 
%%$\displaystyle F(x)=\exp\{-e^{-x}\}$, then
%%\begin{eqnarray}
%%P(Y\leq r|\textbf{x}) &\!\!=\!\!& \exp\{-e^{-(\theta_r+\textbf{x}^T\mbox{\boldmath $\gamma$})}\}
%%\label{eq-3.3.9} \\   \Longleftrightarrow \quad
%%\underbrace{\log\left[-\log P(Y\leq r|\textbf{x})\right]}_{\mbox{log-log link}}
%%&\!\!=\!\!& -(\theta_r+\textbf{x}^T\mbox{\boldmath $\gamma$}), \; r=1,\cdots,q.
%%\label{eq-3.3.10}
%%\end{eqnarray}
%%Model (\ref{eq-3.3.9}) or (\ref{eq-3.3.10}) is referred to as the \textbf{extreme maximal-value distribution
%%model}.
%%\end{itemize}
%%\end{frame}
%%
%%\begin{frame} 
%%Perhaps it is not hard to see there is a connection between the complementary log-log link and the log-log link models above. One can use model (5) to estimate model (4) by:
%%
%%\begin{enumerate}
%%\item Define $\tilde{Y} = k + 1 - Y$, and estimate model (5) with $\tilde{Y}$ (reverse the categories)
%%\item multiple the estimate of (5) with -1 to get estimate for (4). 
%%\end{enumerate}
%%
%%\end{frame}
%%
%%
%%\subsection{\S3.3.2 Extended version of cumulative models}\label{sec:3.3.2}
%%\begin{frame}{S3.3.2 Extended version of cumulative models (1)}
%%%\begin{itemize}
%%%\item
%%%Recall that an ordinal categorical response $Y$ with $k$ categories is determined by a latent variable $U$ as
%%%$$Y\!=\!r \;\Leftrightarrow\; \theta_{r\!-\!1}\!<\!U\!\leq\! \theta_r; \quad r\!=\!1,\!\cdots\!,k;
%%% -\infty\!=\!\theta_0<\theta_1<\!\cdots\! <\theta_k\!=\!+\infty$$
%%%and $U=-\textbf{x}^T\mbox{\boldmath $\gamma$}+\varepsilon$, with $\varepsilon$ having cdf $F$.
%%%\item
%%%Then a \textbf{cumulative model} or \textbf{threshold model} with cdf $F$ is
%%%\begin{equation} 
%%%P(Y\leq r |\textbf{x})=P(U\leq \theta_r)=F(\theta_r+\textbf{x}^T\mbox{\boldmath $\gamma$}).
%%%\label{eq-3.3.3}
%%%\end{equation}
%%
%%\begin{itemize}
%%\item
%%More generally, one may assume a linear form for the thresholds $\theta_1,\cdots, \theta_q$,
%%$$\theta_r=\beta_{r0}+\textbf{w}^T\mbox{\boldmath $\beta$}_r$$
%%where $\textbf{w}$ is a vector explanatory variable and $\mbox{\boldmath $\beta$}_r=(\beta_{r1},
%%\cdots,\beta_{rm})^T$ is a \textit{category-specific} parameter vector.
%%\item 
%%Then an \textbf{extended cumulative model} is obtained:
%%\begin{equation}
%%P(Y\leq r|\textbf{x}, \textbf{w})=F(\beta_{r0}+\textbf{w}^T\mbox{\boldmath $\beta$}_r+
%%\textbf{x}^T\mbox{\boldmath $\gamma$}). \label{eq-3.3.11} 
%%\end{equation}
%%\end{itemize}
%%\end{frame}
%%
%%
%%\begin{frame}{S3.3.2 Extended version of cumulative models (2)}
%%\textbf{Remarks}
%%\begin{enumerate}
%%\item
%%Model (\ref{eq-3.3.11}) becomes unidentifiable if $\textbf{w}_i=\textbf{x}_i$ for each
%%individual $i=1,\cdots,n$; i.e., it will not be possible to separate $\mbox{\boldmath $\gamma$}$ from
%%each $\mbox{\boldmath $\beta$}_r$ if $\textbf{w}_i=\textbf{x}_i$.
%%\item
%%So $\textbf{w}_i$ and $\textbf{x}_i$ must not be the same.
%%\item
%%%$\mbox{\boldmath $\beta$}_r$ can be global, while 
%%$\mbox{\boldmath $\gamma$}$ can be replaced
%%by a category-specific parameter.
%%\item Terminology: 
%%\begin{itemize} \item
%%$\textbf{w}_i$:  \textbf{threshold variables}. \item $\textbf{x}_i$:  
%%\textbf{shift variables}.
%%\end{itemize}
%%\end{enumerate}
%%\end{frame}
%%
%%\subsection{\S3.3.3 Link functions and design matrices for cumulative models}\label{sec:3.3.3}
%%\begin{frame}{\S3.3.3 Cumulative models as GLMs}
%%\begin{itemize}
%%\item
%%For a cumulative GLM, the link function $\mathbf{g}\!=\!(g_1,\cdots, g_q)$,
%%as a vector function of $(\pi_1=P(Y\!=\!1),\cdots, \pi_q=P(Y\!=\!q))$, is given by
%%$$g_r(\pi_1,\cdots,\pi_q)=F^{-1}(\pi_1+\cdots+\pi_r), \quad r=1,\cdots,q,$$
%%where $F^{-1}(\cdot)$ is the inverse function for $F$ that specifies the cumulative model
%%$$\pi_1+\cdots+\pi_r=P(Y\leq r|\textbf{x}, \textbf{w})=F(\beta_{r0}+\textbf{w}^T\mbox{\boldmath $\beta$}_r+
%%\textbf{x}^T\mbox{\boldmath $\gamma$}).$$
%%\item Note that with a link function, one doesn't necessarily have to motivate our models with the latent variable.
%%\end{itemize}
%%\end{frame}
%%
%%
%%\begin{frame}{Different link functions}
%%\begin{itemize}
%%\item
%%$F$ being the $N(0,1)$ cdf gives the \textbf{probit link} 
%%$$g_r(\pi_1,\cdots,\pi_q)=\Phi^{-1}(\pi_1+\cdots+\pi_r), \quad r=1,\cdots,q.$$
%%\item
%%For proportional odds model, the link is the \textbf{cumulative logit}
%%$$g_r(\pi_1,\cdots,\pi_q)=\log\frac{\pi_1+\cdots+\pi_r}{1-(\pi_1+\cdots+\pi_r)}, \quad r=1,\cdots,q.$$
%%\item
%%For proportional hazards model, the link is the \textbf{complement log-log}
%%$$g_r(\pi_1,\cdots,\pi_q)=\log\{-\log(1-(\pi_1+\cdots+\pi_r))\}, \quad r=1,\cdots,q.$$
%%\end{itemize}
%%
%%\end{frame}
%%
%%\begin{frame}{\S3.3.3 Design matrices for cumulative models (1)}
%%\begin{itemize}
%%\item
%%For the simple cumulative model
%%$$P(Y_i\leq r|\textbf{x}_i)=F(\theta_r+\textbf{x}_i^T\mbox{\boldmath $\gamma$}),$$
%%the design matrix $Z_i$ in the linear predictor $\mbox{\boldmath $\eta$}_i=Z_i\mbox{\boldmath $\beta$}$
%%is given by
%%$$Z_i=\left[\begin{array}{ccccc}1 & & & & \textbf{x}_i^T \\
%%& 1 & & &  \textbf{x}_i^T \\ & & \ddots & & \vdots \\
%%& & & 1 & \textbf{x}_i^T\end{array}\right]_{q\times (q+\dim(\textbf{x}_i))}$$
%%and 
%%$$\mbox{\boldmath $\beta$}=(\theta_1,\cdots, \theta_q, \mbox{\boldmath $\gamma$}^T)^T.$$
%%Here $\textbf{x}_i$ should not contain a constant component.
%%\end{itemize}
%%\end{frame}
%%
%%
%%
%%\begin{frame}{\S3.3.3 Design matrices for cumulative models (2)}
%%\begin{itemize}
%%\item
%%For the extended cumulative model
%%$$P(Y_i\leq r|\textbf{x}_i, \textbf{w}_i)=F(\beta_{r0}+\textbf{w}_i^T\mbox{\boldmath $\beta$}_r+
%%\textbf{x}_i^T\mbox{\boldmath $\gamma$}),$$ 
%%the design
%%matrix $Z_i$ in the linear predictor $\mbox{\boldmath $\eta$}_i=Z_i\mbox{\boldmath $\beta$}$
%%is given by
%%$$Z_i=\left[\begin{array}{cccccccc}1 & \textbf{w}_i^T & & & & & & \textbf{x}_i^T \\
%%& & 1 & \textbf{w}_i^T & & & &  \textbf{x}_i^T \\ & & & \ddots & \ddots & & & \vdots \\
%%& & & & & 1 & \textbf{w}_i^T & \textbf{x}_i^T\end{array}\right]_{q\times (q(1+\dim(\textbf{w}_i))+\dim(\textbf{x}_i))}$$
%%and 
%%$$\mbox{\boldmath $\beta$}=(\beta_{10},\mbox{\boldmath $\beta$}_1^T,\cdots, \beta_{q0},
%%\mbox{\boldmath $\beta$}_q^T, \mbox{\boldmath $\gamma$}^T)^T.$$
%%Here neither $\textbf{w}_i$ nor $\textbf{x}_i$ should contain a constant component.
%%\end{itemize}
%%\end{frame}
%%
%%\begin{frame}{\S3.3.3 Alternative formulation/parameterisation in cumulative models (1)}
%%\begin{itemize}
%%\item
%%The threshold parameters $\theta_1,\cdots,\theta_q$ in cumulative models 
%%$P(Y\leq r|\textbf{x})=F(\theta_r+\textbf{x}^T\mbox{\boldmath $\gamma$})$ are restricted by
%%$$\theta_1<\theta_2<\cdots<\theta_q$$
%%\item
%%The corresponding restriction for the extended cumulative model is
%%$$\beta_{10}+\textbf{w}^T\mbox{\boldmath $\beta$}_1<\beta_{20}+\textbf{w}^T\mbox{\boldmath $\beta$}_2
%%<\cdots < \beta_{q0}+\textbf{w}^T\mbox{\boldmath $\beta$}_q$$
%%which is more complex.
%%\item
%%If such constraints are not explicitly accounted for in estimation, the iterative estimation procedure
%%may fail due to fitting inadmissible parameters.
%%\end{itemize}
%%\end{frame}
%%
%%\begin{frame}{\S3.3.3 Alternative formulation/parameterisation in cumulative models (2)}
%%\begin{itemize}
%%\item
%%For the simple cumulative model, the problem can be avoided by using an alternative formulation
%%or reparameterisation:
%%\begin{eqnarray*}
%%\alpha_1=\theta_1,\quad \alpha_r=\log (\theta_r-\theta_{r-1}), \quad r=2,\cdots, q \\
%%\Longleftrightarrow\quad \theta_1=\alpha_1,\quad \theta_r=\theta_1+\sum_{i=2}^r e^{\alpha_i}, \quad
%%r=2,\cdots,q.
%%\end{eqnarray*}
%%The new parameters $\alpha_1,\cdots,\alpha_q$ are unconstrained.
%%\item
%%The cumulative model under the new formulation has a linear structure of the following form
%%\begin{eqnarray*}
%%F^{-1}\left(P(Y=1|\textbf{x})\right)=\alpha_1+\textbf{x}^T\mbox{\boldmath $\gamma$\hspace{45pt}} \\
%%\log\left[F^{-1}\left(P(Y\leq r|\textbf{x})\right)-F^{-1}\left(P(Y\leq r-1|\textbf{x})\right)\right]=\alpha_r,
%%\quad r=2,\cdots,q.
%%\end{eqnarray*}
%%\item
%%The link function corresponding to the new formulation can be determined accordingly.
%%\end{itemize}
%%\end{frame}
%%
%%\begin{frame}{\S3.3.3 Alternative formulation/parameterisation in cumulative models (3)}
%%\begin{itemize}
%%\item
%%For the special case of the \textit{cumulative logistic model}, one obtains the following link function
%%(\textbf{log-logit link}) {\footnotesize
%%\begin{eqnarray*}
%%g_1(\pi_1,\cdots,\pi_q) &=& \log\left(\frac{\pi_1}{1-\pi_1}\right) \\
%%g_r(\pi_1,\cdots,\pi_q) &=& \log\left[\log\left\{\frac{\pi_1+\cdots+\pi_r}{1-\pi_1-\cdots-\pi_r}\right\} -
%%\log\left\{\frac{\pi_1+\cdots+\pi_{r-1}}{1-\pi_1-\cdots-\pi_{r-1}}\right\}\right], \\ & & \quad q=2,\cdots,q.
%%\end{eqnarray*}}
%%\item
%%If this alternative link function is used, the design matrix $Z_i$ has to be adapted. Now {\footnotesize
%%$$Z_i=\left[\begin{array}{ccccc}1 & & & & \textbf{x}_i^T \\
%%& 1 & & &  \textbf{0} \\ & & \ddots & & \vdots \\
%%& & & 1 & \textbf{0}\end{array}\right]_{q\times (q+\dim(\textbf{x}_i))}$$}
%%and the parameter vector becomes
%%$\mbox{\boldmath $\beta$}=(\alpha_1,\cdots, \alpha_q, \mbox{\boldmath $\gamma$}^T)^T.$
%%
%%Here $\textbf{x}_i$ should not contain a constant component.
%%\end{itemize}
%%\end{frame}
%%
%\subsection{\S3.3.3 Link functions and design matrices for cumulative models}\label{sec:3.3.3}
%\begin{frame}{\S3.3.3 Computing packages and an example}
%\begin{itemize}
%\item
%The \texttt{polr()} function in the R \textbf{MASS} package can be used to fit cumulative models
%(but not extended cumulative models).
%\item
%More generally, the \textbf{ordinal} package in R can be used for fitting ordinal regression models.
%\end{itemize}
%
%\textbf{Example 3.4} (Example 2 in Chapter 1). Breathing test results.
%\begin{itemize}
%\item
%Response variable: \texttt{BTR}, ordinal with $k=3$ levels (1 = Normal, 2 = Border, 3 = Abnormal ).
%\item
%\texttt{Age} (2 levels) and \texttt{Smoking} status (3 levels) are predictors.
%\item
%Aim: see how \texttt{Age} and \texttt{Smoking} are associated with \texttt{BTR}.
%\item
%Command \texttt{polr} in R \texttt{MASS} package is used in the analysis.
%\end{itemize}
%\end{frame}
%
%\begin{frame}[fragile] \frametitle{\S3.3.3 Breathing testing example (1)}
%\textbf{Usage of} \texttt{polr}: {\footnotesize
%\begin{verbatim}
%polr(formula, data, weights, start, ..., subset, na.action,
%     contrasts = NULL, Hess = FALSE, model = TRUE,
%     method = c("logistic", "probit", "loglog", "cloglog", "cauchit"))
%\end{verbatim}}
%Note the estimates of $\mbox{\boldmath $\gamma$}$ are those from \texttt{polr} multiplied by $-1$
%due to the model used in \texttt{polr} being $P(Y\leq r|\textbf{x})=F(\theta_r-\textbf{x}^T\mbox{\boldmath
%$\gamma$})$.
%\begin{tiny}
%\begin{Schunk}
%\begin{Sinput}
%> library(MASS); BTR.dat <- read.csv("D:/MAST90084/Example3-2BTR.csv"); BTR.dat
%\end{Sinput} 
%\begin{Soutput}
%       BTR    Age  Smoking Freq
%1  1Normal    <40   1Never  577
%2  2Border    <40   1Never   27
%3  3Abnorm    <40   1Never    7
%4  1Normal    <40  2Former  192
%5  2Border    <40  2Former   20
%6  3Abnorm    <40  2Former    3
%7  1Normal    <40 3Current  682
%8  2Border    <40 3Current   46
%9  3Abnorm    <40 3Current   11
%10 1Normal 40to59   1Never  164
%11 2Border 40to59   1Never    4
%12 3Abnorm 40to59   1Never    0
%13 1Normal 40to59  2Former  145
%14 2Border 40to59  2Former   15
%15 3Abnorm 40to59  2Former    7
%16 1Normal 40to59 3Current  245
%17 2Border 40to59 3Current   47
%18 3Abnorm 40to59 3Current   27
%\end{Soutput}
%\end{Schunk}
%\end{tiny}
%\end{frame}
%
%\begin{frame}[fragile] \frametitle{\S3.3.3 Breathing testing example (2)}
%\begin{scriptsize}
%\begin{Schunk}
%\begin{Sinput}
%> is.factor(BTR.dat$BTR)    #TRUE;     > is.ordered(BTR.dat$BTR)    #FALSE
%> is.factor(BTR.dat$Age)    #TRUE;     > is.factor(BTR.dat$Smoking)  #TRUE
%> is.factor(BTR.dat$Freq) #FALSE;     
%> BTR.dat$BTR=as.ordered(BTR.dat$BTR)   #Set BTR as an ordinal factor
%> help(polr)
%> BTR.logit=polr(BTR~Age+Smoking+Age:Smoking, data=BTR.dat, weights=Freq, 
%                  Hess=T,method="logistic")
%> summary(BTR.logit)
%\end{Sinput} 
%\begin{Soutput}
%Coefficients:
%                            Value Std. Error t value
%Age40to59                 -0.8857     0.5359  -1.653
%Smoking2Former             0.6973     0.2822   2.471
%Smoking3Current            0.3472     0.2238   1.551
%Age40to59:Smoking2Former   1.1458     0.6228   1.840
%Age40to59:Smoking3Current  2.2007     0.5689   3.868
%
%Intercepts:
%                Value   Std. Error t value
%1Normal|2Border  2.8334  0.1764    16.0603
%2Border|3Abnorm  4.3078  0.2130    20.2210
%
%Residual Deviance: 1564.968 
%AIC: 1578.968 
%\end{Soutput}
%\end{Schunk}
%\end{scriptsize}
%\end{frame}
%
%\begin{frame}[fragile] \frametitle{\S3.3.3 Breathing testing example (3)}
%\begin{scriptsize}
%\begin{Schunk}
%\begin{Sinput}
%> BTR.logit1=polr(BTR~Age+Smoking, data=BTR.dat, weights=Freq, Hess=T,method="logistic")
%> summary(BTR.logit1)
%\end{Sinput} 
%\begin{Soutput}
%Coefficients:
%                 Value Std. Error t value
%Age40to59       0.7772     0.1482   5.243
%Smoking2Former  0.7815     0.2333   3.350
%Smoking3Current 0.9607     0.1918   5.010
%
%Intercepts:
%                Value   Std. Error t value
%1Normal|2Border  3.1927  0.1749    18.2544
%2Border|3Abnorm  4.6543  0.2125    21.9076
%
%Residual Deviance: 1589.744 
%AIC: 1599.744 
%\end{Soutput}
%\begin{Sinput}
%> anova(BTR.logit, BTR.logit1)
%\end{Sinput} 
%\begin{Soutput}
%Likelihood ratio tests of ordinal regression models
%
%Response: BTR
%                        Model Resid. df Resid. Dev   Test    Df LR stat.   Pr(Chi)
%1               Age + Smoking      2214   1589.744                                   
%2 Age + Smoking + Age:Smoking      2212   1564.968 1 vs 2     2 24.77585 4.168e-06
%\end{Soutput}
%\end{Schunk}
%\end{scriptsize}
%\end{frame}
%
%\begin{frame}[fragile] \frametitle{\S3.3.3 Breathing testing example (4)}
%\begin{tiny}
%\begin{Schunk}
%\begin{Sinput}
%> BTR.logit$fitted
%\end{Sinput} 
%\begin{Soutput}
%     1Normal    2Border     3Abnorm
%1  0.9444565 0.04225961 0.013283873
%2  0.9444565 0.04225961 0.013283873
%3  0.9444565 0.04225961 0.013283873
%4  0.8943685 0.07930628 0.026325236
%5  0.8943685 0.07930628 0.026325236
%6  0.8943685 0.07930628 0.026325236
%7  0.9231700 0.05813458 0.018695369
%8  0.9231700 0.05813458 0.018695369
%9  0.9231700 0.05813458 0.018695369
%10 0.9763207 0.01815786 0.005521451
%11 0.9763207 0.01815786 0.005521451
%12 0.9763207 0.01815786 0.005521451
%13 0.8671630 0.09895796 0.033879045
%14 0.8671630 0.09895796 0.033879045
%15 0.8671630 0.09895796 0.033879045
%16 0.7633793 0.17036521 0.066255462
%17 0.7633793 0.17036521 0.066255462
%18 0.7633793 0.17036521 0.066255462
%\end{Soutput}
%\begin{Sinput}
%> BTR.logit$residual  #does not exist.
%> BTR.logit$lp  #linear predictor values
%\end{Sinput} 
%\begin{Soutput}
%         1          2          3          4          5          6          7          8          9 
% 0.0000000  0.0000000  0.0000000  0.6972823  0.6972823  0.6972823  0.3472245  0.3472245  0.3472245 
%        10         11         12         13         14         15         16         17         18 
%-0.8857462 -0.8857462 -0.8857462  0.9573393  0.9573393  0.9573393  1.6621466  1.6621466  1.6621466 
%\end{Soutput}
%\begin{Sinput}
%> attributes(BTR.logit)
%\end{Sinput} 
%\begin{Soutput}
% [1] "coefficients"  "zeta"          "deviance"      "fitted.values" "lev"           "terms"        
% [7] "df.residual"   "edf"           "n"             "nobs"          "call"          "method"       
%[13] "convergence"   "niter"         "lp"            "Hessian"       "model"         "contrasts"    
%[19] "xlevels"      
%$class
%[1] "polr"
%\end{Soutput}
%\end{Schunk}
%\end{tiny}
%\end{frame}
%
%\begin{frame}{\S3.3.3 Breathing testing example (5)}
%\begin{itemize}
%\item 
%On the previous 3 slides we fit two cumulative logistic models to \texttt{BTR} using 
%\texttt{Age} and \texttt{Smoking}: the first (\texttt{BTR.logit} includes the \texttt{Age:Smokings} 
%interaction-effect terms as well as the main-effect terms; while the second (\texttt{BTR.logit1}
%includes main-effect terms only.
%\item
%Dummy coding (i.e. \texttt{contr.treatment} coding in R) is used for representing \texttt{Age} and \texttt{Smoking}:
%{\footnotesize
%\begin{eqnarray*}
%\texttt{Age40to59}&=& \left\{\begin{array}{cc} 1, & \mbox{if age is 40 $\sim$ 59} \\ 0 & \mbox{if age < 40}
%\end{array}\right. \\
%\texttt{Smoking2Former}&=&\left\{\begin{array}{cc} 1, & \mbox{former smoker} \\ 0 & \mbox{else}
%\end{array}\right.,\\
%\texttt{Smoking3Current}&=&\left\{\begin{array}{cc} 1, & \mbox{current smoker} \\ 0 & \mbox{else}
%\end{array}\right.
%\end{eqnarray*}}
%\end{itemize}
%\end{frame}
%
%\begin{frame}{\S3.3.3 Breathing testing example (6)}
%\begin{itemize}
%\item
%The first model can be written as {\scriptsize
%\begin{eqnarray*}
%\log\frac{P(\texttt{BTR}=\texttt{Normal})}{P(\texttt{BTR}>\texttt{Normal})}\!\!&\!\!\!\!=\!\!\!&\!\!\!\theta_1
%+ \gamma_1\texttt{Age40to59}+\gamma_2\texttt{Smoking2Former}+\gamma_3\texttt{Smoking3Current} \\
%\!\!\!\!&\!\!\!\! &\!\!\!\!\!\!\!\!\!\!\!\!\!\!\!\!+ \gamma_4\texttt{Age40to59}\cdot \texttt{Smoking2Former}
%+ \gamma_5\texttt{Age40to59}\cdot \texttt{Smoking3Current} \\
%\log\frac{P(\texttt{BTR}\leq \texttt{Bordering})}{P(\texttt{BTR}>\texttt{Bordering})}\!\!&\!\!\!\!=\!\!\!&\!\!\!\theta_2
%+ \gamma_1\texttt{Age40to59}+\gamma_2\texttt{Smoking2Former}+\gamma_3\texttt{Smoking3Current} \\
%\!\!\!\!&\!\!\!\! &\!\!\!\!\!\!\!\!\!\!\!\!\!\!\!\! + \gamma_4\texttt{Age40to59}\cdot \texttt{Smoking2Former}
%+ \gamma_5\texttt{Age40to59}\cdot \texttt{Smoking3Current} 
%\end{eqnarray*}}
%\item
%MLEs of the parameters in the first model are {\scriptsize
%$$\hat{\theta}_1=2.833, \hat{\theta}_2=4.308, \hat{\gamma}_1=0.886, \hat{\gamma}_2=-0.697,
%\hat{\gamma}_3=-0.347, \hat{\gamma}_4=-1.146, \hat{\gamma}_5=-2.201$$}
%with their standard errors given in the R output.
%\item
%Estimated values of the $\gamma$ parameters can be interpreted as log cumulative odds ratios.
%\textbf{For example}, the odds of having a normal test result (or normal or bordering result)
%for a worker aged 40 to 59 who never smoked 
%is estimated to be $e^{-\hat{\gamma}_2-\hat{\gamma}_4}=e^{0.697+1.146}=6.32$ times of that for a worker
%aged 40 to 59 who is a former worker.
%\end{itemize}
%\end{frame}
%
%\begin{frame}{\S3.3.3 Breathing testing example (7)}
%\begin{itemize}
%\item
%The second model can be written as {\scriptsize
%\begin{eqnarray*}
%\log\frac{P(\texttt{BTR}=\texttt{Normal})}{P(\texttt{BTR}>\texttt{Normal})}\!\!&\!\!\!\!=\!\!\!&\!\!\!\theta_1
%+ \gamma_1\texttt{Age40to59}+\gamma_2\texttt{Smoking2Former}+\gamma_3\texttt{Smoking3Current} \\
%\log\frac{P(\texttt{BTR}\leq \texttt{Bordering})}{P(\texttt{BTR}>\texttt{Bordering})}\!\!&\!\!\!\!=\!\!\!&\!\!\!\theta_2
%+ \gamma_1\texttt{Age40to59}+\gamma_2\texttt{Smoking2Former}+\gamma_3\texttt{Smoking3Current}
%\end{eqnarray*}}
%\item
%MLEs of the parameters in the first model are {\scriptsize
%$$\hat{\theta}_1=3.193, \;\hat{\theta}_2=4.654, \; \hat{\gamma}_1=-0.777, \;\hat{\gamma}_2=-0.782, \;
%\hat{\gamma}_3=-0.961$$}
%with their standard errors given in the R output.
%\item
%``$H_0: \gamma_4=\gamma_5=0$ vs. $H_1: \mbox{not}\; H_0$" is the hypothesis about the 
%\texttt{Age:Smoking} interaction effects on \texttt{BTR}.  
%\item
%A $\chi^2$ analysis of deviance test comparing models \texttt{BTR.logit} and texttt{BTR.logit1} gives 
%an LR statistic of 22.776 and $p$-value of $4.168\times 10^{-6}$. Therefore, there is very strong evidence to reject $H_0$.
%We conclude that \texttt{Age} and \texttt{Smoking} have very strong interaction effects on \texttt{BTR}.
%\end{itemize}
%\end{frame}
%
%\begin{frame}{\S3.3.3 Breathing testing example (8)}
%\begin{itemize}
%\item
%Linear predictor of the first model is {\scriptsize
%\begin{eqnarray*}
%\eta_1 &=& \theta_1
%+ \gamma_1\texttt{Age40to59}+\gamma_2\texttt{Smoking2Former}+\gamma_3\texttt{Smoking3Current} \\
%& & + \gamma_4\texttt{Age40to59}\cdot \texttt{Smoking2Former}
%+ \gamma_5\texttt{Age40to59}\cdot \texttt{Smoking3Current} \\
%\eta_2&=& \theta_2
%+ \gamma_1\texttt{Age40to59}+\gamma_2\texttt{Smoking2Former}+\gamma_3\texttt{Smoking3Current} \\
%& & + \gamma_4\texttt{Age40to59}\cdot \texttt{Smoking2Former}
%+ \gamma_5\texttt{Age40to59}\cdot \texttt{Smoking3Current} 
%\end{eqnarray*}}
%\item
%However, the linear predictor returned from \texttt{polr} is just $\textbf{x}^T\mbox{\boldmath $\gamma$}$, i.e.
%{\scriptsize
%\begin{eqnarray*}
%\lefteqn{ \gamma_1\texttt{Age40to59}+\gamma_2\texttt{Smoking2Former}+\gamma_3\texttt{Smoking3Current}} & \\
%& &  + \gamma_4\texttt{Age40to59}\cdot \texttt{Smoking2Former}
%+ \gamma_5\texttt{Age40to59}\cdot \texttt{Smoking3Current}
%\end{eqnarray*}}
%thus has just 6 different values knowing that there are 6 combinations of \texttt{Age} and \texttt{Smoking}.
%\item
%Knowing these tricks, there should be no difficulty to reveal how the fitted values returned by
%\texttt{BTR.logit\$fitted} are obtained.
%\end{itemize}
%\end{frame}
%
%\begin{frame}{\S3.3.3 Breathing testing example (9)}
%\begin{itemize}
%\item
%Effect coding (i.e. \texttt{contr.sum} coding in R, $-1$ for last category) 
%may also be used for representing \texttt{Age} and \texttt{Smoking}:
%{\scriptsize
%\begin{eqnarray*}
%\texttt{Age1}&=& \left\{\begin{array}{cc} 1, & \mbox{if age < 40} \\ -1, & \mbox{if age is 40 $\sim$ 59}
%\end{array}\right. \\
%\texttt{Smoking1}&=&\left\{\begin{array}{cc} 1, & \mbox{never smoked} \\ 0, & \mbox{former smoker} \\ -1, & 
%\mbox{current smoker} \end{array}\right., \quad
%\texttt{Smoking2}\;=\; \left\{\begin{array}{cc} 0, & \mbox{never smoked} \\ 1, & \mbox{former smoker} \\ 
%-1, & \mbox{current smoker}
%\end{array}\right.
%\end{eqnarray*}}
%\item
%With this coding, the first model can be written as {\scriptsize
%\begin{eqnarray*}
%\log\frac{P(\texttt{BTR}=\texttt{Normal})}{P(\texttt{BTR}>\texttt{Normal})}&=&\theta_1
%+ \gamma_1\texttt{Age1}+\gamma_2\texttt{Smoking1}+\gamma_3\texttt{Smoking2} \\
%&& + \gamma_4\texttt{Age1}\cdot \texttt{Smoking1}
%+ \gamma_5\texttt{Age1}\cdot \texttt{Smoking2} \\
%\log\frac{P(\texttt{BTR}\leq \texttt{Bordering})}{P(\texttt{BTR}>\texttt{Bordering})}&=&\theta_2
%+ \gamma_1\texttt{Age1}+\gamma_2\texttt{Smoking1}+\gamma_3\texttt{Smoking2} \\
%&& + \gamma_4\texttt{Age1}\cdot \texttt{Smoking1}
%+ \gamma_5\texttt{Age1}\cdot \texttt{Smoking2} 
%\end{eqnarray*}}
%\item
%Although the parameter estimates will be different now, other results from the model will be the
%same as that using dummy coding.
%\end{itemize}
%\end{frame}
%
%
%
%\begin{frame} {Coding scheme revisited (2.1 in F \& T)}
% 
%
%
%\end{frame}
%
%\begin{frame}[fragile] \frametitle{\S3.3.3 Breathing testing example (10)}
%\begin{tiny}
%\begin{Schunk}
%\begin{Sinput}
%> options(contrasts = c("contr.sum", "contr.poly"))
%> BTR.logit=polr(BTR~Age+Smoking+Age:Smoking, data=BTR.dat, weights=Freq, Hess=T,method="logistic")
%> summary(BTR.logit)
%\end{Sinput} 
%\begin{Soutput}
%Coefficients:
%                 Value Std. Error t value
%Age1          -0.11487     0.1086 -1.0580
%Smoking1      -0.90596     0.1890 -4.7933
%Smoking2       0.36428     0.1421  2.5644
%Age1:Smoking1  0.55777     0.1890  2.9512
%Age1:Smoking2 -0.01517     0.1421 -0.1068
%
%Intercepts:
%                Value   Std. Error t value
%1Normal|2Border  2.3704  0.1086    21.8209
%2Border|3Abnorm  3.8447  0.1599    24.0486
%
%Residual Deviance: 1564.968 
%AIC: 1578.968 
%\end{Soutput}
%\begin{Sinput}
%> BTR.logit1=polr(BTR~Age+Smoking, data=BTR.dat, weights=Freq, Hess=T,method="logistic")
%> anova(BTR.logit, BTR.logit1)
%\end{Sinput} 
%\begin{Soutput}
%Likelihood ratio tests of ordinal regression models
%
%Response: BTR
%                        Model Resid. df Resid. Dev   Test    Df LR stat.      Pr(Chi)
%1               Age + Smoking      2214   1589.744                                   
%2 Age + Smoking + Age:Smoking      2212   1564.968 1 vs 2     2 24.77585 4.168618e-06
%\end{Soutput}
%\end{Schunk}
%\end{tiny}
%\end{frame}
%
%\begin{frame}[fragile] \frametitle{\S3.3.3 Breathing testing example (11)}
%\begin{tiny}
%\begin{Schunk}
%\begin{Sinput}
%> attributes(BTR.logit)
%\end{Sinput} 
%\begin{Soutput}
%$names
% [1] "coefficients"  "zeta"          "deviance"      "fitted.values" "lev"           "terms"         "df.residual"  
% [8] "edf"           "n"             "nobs"          "call"          "method"        "convergence"   "niter"        
%[15] "lp"            "Hessian"       "model"         "contrasts"     "xlevels"      
%\end{Soutput}
%\begin{Sinput}
%> BTR.logit$fitted
%\end{Sinput} 
%\begin{Soutput}
%     1Normal    2Border     3Abnorm
%1  0.9444565 0.04225911 0.013284361
%2  0.9444565 0.04225911 0.013284361
%3  0.9444565 0.04225911 0.013284361
%4  0.8943660 0.07930714 0.026326874
%5  0.8943660 0.07930714 0.026326874
%6  0.8943660 0.07930714 0.026326874
%7  0.9231672 0.05813603 0.018696801
%8  0.9231672 0.05813603 0.018696801
%9  0.9231672 0.05813603 0.018696801
%10 0.9763222 0.01815653 0.005521309
%11 0.9763222 0.01815653 0.005521309
%12 0.9763222 0.01815653 0.005521309
%13 0.8671585 0.09895993 0.033881540
%14 0.8671585 0.09895993 0.033881540
%15 0.8671585 0.09895993 0.033881540
%16 0.7633676 0.17037058 0.066261785
%17 0.7633676 0.17037058 0.066261785
%18 0.7633676 0.17037058 0.066261785
%\end{Soutput}
%\begin{Sinput}
%> BTR.logit$lp    #linear predictor values
%\end{Sinput} 
%\begin{Soutput}
%         1          2          3          4          5          6          7          8          9         10 
%-0.4630592 -0.4630592 -0.4630592  0.2342498  0.2342498  0.2342498 -0.1157938 -0.1157938 -0.1157938 -1.3488684 
%        11         12         13         14         15         16         17         18 
%-1.3488684 -1.3488684  0.4943191  0.4943191  0.4943191  1.1991525  1.1991525  1.1991525
%\end{Soutput}
%\end{Schunk}
%\end{tiny}
%\end{frame}
%
%\begin{frame}[fragile] \frametitle{\S3.3.3 Breathing testing example (12)}
%We also fit a grouped Cox or proportional hazards model.
%\begin{scriptsize}
%\begin{Schunk}
%\begin{Sinput}
%> BTR.ph=polr(BTR~Age+Smoking+Age:Smoking,data=BTR.dat, weights=Freq, 
%                      Hess=T, method="cloglog")
%> summary(BTR.ph)
%\end{Sinput} 
%\begin{Soutput}
%Coefficients:
%                            Value Std. Error t value
%Age40to59                 -0.2865    0.14264  -2.008
%Smoking2Former             0.2125    0.10054   2.114
%Smoking3Current            0.1088    0.07301   1.490
%Age40to59:Smoking2Former   0.4333    0.18941   2.288
%Age40to59:Smoking3Current  0.8379    0.16357   5.122
%
%Intercepts:
%                Value   Std. Error t value
%1Normal|2Border  1.0477  0.0558    18.7612
%2Border|3Abnorm  1.5528  0.0629    24.6916
%
%Residual Deviance: 1559.949 
%AIC: 1573.949 
%\end{Soutput}
%\end{Schunk}
%\end{scriptsize}
%\end{frame}
%
%\begin{frame}[fragile] \frametitle{\S3.3.3 Breathing testing example (13)}
%We further fit a cumulative model with the extreme maximal-value distribution.
%\begin{scriptsize}
%\begin{Schunk}
%\begin{Sinput}
%> BTR.extm=polr(BTR~Age+Smoking+Age:Smoking,data=BTR.dat,weights=Freq, 
%                              Hess=T,method="loglog")
%> summary(BTR.extm)
%\end{Sinput} 
%\begin{Soutput}
%Coefficients:
%                            Value Std. Error t value
%Age40to59                 -0.8672     0.5286  -1.640
%Smoking2Former             0.6755     0.2700   2.501
%Smoking3Current            0.3368     0.2167   1.554
%Age40to59:Smoking2Former   1.1003     0.6070   1.813
%Age40to59:Smoking3Current  2.0724     0.5573   3.719
%
%Intercepts:
%                Value   Std. Error t value
%1Normal|2Border  2.8614  0.1715    16.6838
%2Border|3Abnorm  4.2761  0.2080    20.5605
%
%Residual Deviance: 1566.335 
%AIC: 1580.335 
%\end{Soutput}
%\end{Schunk}
%\end{scriptsize}
%\end{frame}
%
%\begin{frame}[fragile] \frametitle{\S3.3.3 Breathing testing example (14)}
%We fit a grouped Cox or proportional hazards model using the effect coding.
%\begin{scriptsize}
%\begin{Schunk}
%\begin{Sinput}
%> options(contrasts = c("contr.sum", "contr.poly"))
%> BTR.ph=polr(BTR~Age+Smoking+Age:Smoking,data=BTR.dat, weights=Freq, 
%                           Hess=T, method="cloglog")
%> summary(BTR.ph)
%\end{Sinput} 
%\begin{Soutput}
%Coefficients:
%                 Value Std. Error  t value
%Age1          -0.06862    0.03421 -2.00575
%Smoking1      -0.31898    0.05366 -5.94408
%Smoking2       0.11022    0.04961  2.22168
%Age1:Smoking1  0.21188    0.05361  3.95205
%Age1:Smoking2 -0.00481    0.04960 -0.09699
%
%Intercepts:
%                Value   Std. Error t value
%1Normal|2Border  0.8720  0.0353    24.7072
%2Border|3Abnorm  1.3771  0.0441    31.2213
%
%Residual Deviance: 1559.949 
%AIC: 1573.949 
%\end{Soutput}
%\end{Schunk}
%\end{scriptsize}
%\end{frame}
%
%\begin{frame}[fragile] \frametitle{\S3.3.3 Breathing testing example (15)}
%We finally fit a cumulative model with the extreme maximal-value distribution using the effect coding.
%\begin{scriptsize}
%\begin{Schunk}
%\begin{Sinput}
%> options(contrasts = c("contr.sum", "contr.poly"))
%> BTR.extm=polr(BTR~Age+Smoking+Age:Smoking,data=BTR.dat,weights=Freq, 
%                              Hess=T,method="loglog")
%> summary(BTR.extm)
%\end{Sinput} 
%\begin{Soutput}
%Coefficients:
%                 Value Std. Error t value
%Age1          -0.09523     0.1053  -0.904
%Smoking1      -0.86618     0.1854  -4.671
%Smoking2       0.35943     0.1361   2.642
%Age1:Smoking1  0.52877     0.1854   2.852
%Age1:Smoking2 -0.02136     0.1361  -0.157
%
%Intercepts:
%                Value   Std. Error t value
%1Normal|2Border  2.4288  0.1054    23.0536
%2Border|3Abnorm  3.8434  0.1573    24.4373
%
%Residual Deviance: 1566.335 
%AIC: 1580.335 
%\end{Soutput}
%\end{Schunk}
%\end{scriptsize}
%\end{frame}
%
\subsection{\S3.3.4 Sequential models}\label{sec:3.3.4}


\begin{frame} {Sequential models for ordinal responses}


\begin{itemize}
\item {\bf Cumulative models}: As MGLMs,  the (multivariate) link function essentially acts on the cumulative probabilities $P(Y \leq r | {\bf x}) =  \pi_1 + \dots + \pi_r $ and has the form
\[
g_r( \pi_1, \dots, \pi_q)  = F^{-1} ( \pi_1 + \dots + \pi_r )\text{ for all } r = 1, \dots, \underbrace{q}_{k-1},
\]
for some CDF $F$. 

\ \
\item 
{\bf Sequential models}: another class of MGLMs where the link  has the  form
\begin{multline*}
g_r( \pi_1, \dots, \pi_q)  =  F^{-1} \left(\frac{\pi_r}{ 1 - \pi_1 - \cdots -  \pi_{r -1 }  } \right) = F^{-1}\bigg(\frac{P(Y = r)}{P(Y \geq r)} \bigg), \\
\text{ for }
  r = 1, \dots q,
\end{multline*}
 acting on the discrete hazards $P(Y = r)/P(Y \geq r)$.
\end{itemize}
\end{frame}

\begin{frame}{\S3.3.4 Motivating example for sequential models}
%\begin{itemize}
%
%\end{itemize}

\noindent \textbf{Example 3.6. (Tonsil size)} Children have been classified according to their tonsil size ($Y$) 
and whether or not they are carriers of Streptococcus pyogenes ($X$). 
\begin{table} [htbp]
\caption{Tonsil size and Streptococcus pyogenes (Holmes \& Williams, 1954)} \label{tab3.5}
\begin{center} {\small
\begin{tabular}{|lccc|} \hline
 & $\begin{array}{c} \mbox{Present but} \\ \mbox{not enlarged}\end{array}$ & Enlarged & 
 $\begin{array}{c} \mbox{Greatly} \\ \mbox{enlarged}\end{array}$ \\ \hline
Carriers & 19 & 29 & 24 \\
Noncarriers & 497 & 560 & 269 \\ \hline
\end{tabular}}
\end{center}
\end{table}
\begin{itemize}
\item
{\scriptsize Assume that the tonsil size always starts in the normal state ``present
but not enlarged" (\textit{category 1}). If the tonsils grow abnormally, they may
become ``enlarged" (\textit{category 2}); if the process does not stop, tonsils may
become ``greatly enlarged" (\textit{category 3}). But in order to get greatly enlarged
tonsils, they first have to be enlarged for the duration of the intermediate
state ``enlarged."}
\end{itemize}
\end{frame}


\begin{frame}{\S3.3.4 Motivation for sequential models}
\begin{wideitemize}
\item
The response categories are ordered and  can be reached only successively, based on  a sequential mechanism with $q = k - 1$ latent variables $U_r$. 

\item
For $r=1,\cdots,q$, let
a  latent  $U_r$ has the linear form
$$U_r=-\textbf{x}^T\mbox{\boldmath $\gamma$}+\varepsilon_r, \quad\mbox{where $\varepsilon_r$
has cdf $F$.}$$


\item (cf. We only used \underline{one} latent variable to motivate {\bf cumulative models}).
%



% in (\ref{eq-3.3.17}).
\end{wideitemize}
\end{frame}

\begin{frame}
\begin{wideitemize}
\item The value of
$Y$ is determined as, for each $r = 1, \dots, q$:
\[
Y = r \iff  U_1 > \theta_1, \dots, U_{r-1} > \theta_{r-1} \text{ and } U_r \leq \theta_r.
\]
Else, $Y = k = q+1$.


%\begin{eqnarray}
%Y=1 & \Leftrightarrow & U_1\leq\theta_1 \notag \\
%Y=2 \;\mbox{given $Y\geq 2$} & \Leftrightarrow & U_2\leq\theta_2; \quad\mbox{and in general} \nonumber \\
%Y=r \;\mbox{given $Y\geq r$} & \Leftrightarrow & U_r\leq\theta_r; \quad\mbox{or equivalently} \nonumber \\
%Y>r \;\mbox{given $Y\geq r$} & \Leftrightarrow & U_r> \theta_r, \quad r=1,\cdots, q \label{eq-3.3.16}
%\end{eqnarray}
%
%\item In general,

\begin{enumerate}
\item So, if $U_1 \leq \theta_1$, $Y$ determined to be $1$, else $Y \geq 2$;
\item From there,  if $U_2 \leq \theta_2$, $Y$ determined to be $2$;
\item   The process goes on \emph{sequentially } until $Y$ is determined. 
\end{enumerate}


\item It  models the transition from category $r$ to category $r+1$, 
given that category $r$ is reached.
 

\item
\textcolor{red}{NOTE}: No order restriction on $\theta_1, \cdots, \theta_q$.
\end{wideitemize}

\end{frame}

%
%\begin{frame}
%\begin{wideitemize}
%
%\item
%The main difference to the threshold approach used in the cumulative model is the conditional modelling 
%of the transitions. The sequential mechanism assumes a binary decision in each step.
%\end{wideitemize}
%\end{frame}



% new frame
\begin{frame}
\begin{wideitemize}
\item
%The sequential response mechanism (\ref{eq-3.3.16}) combined with 

The linear form of the latent variable $U_r = - {\bf x}^T {\boldsymbol \gamma} + \epsilon_r$ for $\epsilon_r \sim F$
%$$U_r=-\textbf{x}^T\mbox{\boldmath $\gamma$}+\varepsilon_r, \quad\mbox{where $\varepsilon_r$
%has cdf $F$,}$$
then leads to the \textbf{sequential model}
\begin{equation*}
\underbrace{P(Y=r|Y\geq r, \textbf{x})}_{\text{discrete hazard, } = \frac{\pi_r}{1 - \pi_1 - \dots - \pi_{r-1}}}= 
P(U_r\leq \theta_r)=F(\theta_r+\textbf{x}^T\mbox{\boldmath $\gamma$}), 
\label{eq-3.3.17}
\end{equation*}
for $r=1,\cdots,k$, where $\theta_k\equiv +\infty$, i.e. the linear predictor $\theta_r+\textbf{x}^T\mbox{\boldmath $\gamma$}$ models the discrete hazard  via the CDF $F(\cdot)$.


\item \textcolor{red}{Discrete hazard}: the instantaneous chance of dying at ``time" $r$ conditional on living to ``time" $r$.
\item Inverting $F$  gets the  link function components $g_1, \dots, g_q$. 
%\item
%Model (\ref{eq-3.3.17}) is also called the \textbf{continuation ratio model}.
\end{wideitemize}
\end{frame}

\begin{frame}
\begin{itemize}
\item The components of the vector-valued response function can be explicitly recovered as, for each $r = 1, \dots, k$,
{\footnotesize
\begin{eqnarray*}
\lefteqn{\pi_r= P(Y=r|\textbf{x})=P(Y=r|Y\geq r, \textbf{x})\cdot P(Y\geq r |\textbf{x})} & \\
& & =F(\theta_r+\textbf{x}^T\mbox{\boldmath $\gamma$}) \cdot P(Y>r-1|Y\geq r-1, \textbf{x}) \cdot \\
& & \quad\quad \quad P(Y>r-2|Y\geq r-2, \textbf{x})\cdots P(Y>1|Y\geq 1, \textbf{x}) \\
&&=F(\theta_r+\textbf{x}^T\mbox{\boldmath $\gamma$})[1-F(\theta_{r-1}+\textbf{x}^T\mbox{\boldmath $\gamma$})]
\cdots [1-F(\theta_1+\textbf{x}^T\mbox{\boldmath $\gamma$})] \\
&& = F(\theta_r+\textbf{x}^T\mbox{\boldmath $\gamma$})\prod_{i=1}^{r-1}
[1-F(\theta_i+\textbf{x}^T\mbox{\boldmath $\gamma$})].
%; \quad 
%\mbox{Note $\prod_{i=1}^0 [\cdot]=0$}
\end{eqnarray*}
}
\item (We have defined $\prod_{i=1}^0 [\cdot]=1$)
\end{itemize}

\end{frame}

\begin{frame}{\S3.3.4 Some examples}
%\begin{itemize}
%
%\item

\begin{enumerate}
\item
A logistic $F(x)=(1+e^{-x})^{-1}$ gives the \textbf{sequential logit model}:
$$P(Y=r|Y\geq r, \textbf{x})=\frac{\exp(\theta_r+\textbf{x}^T\mbox{\boldmath $\gamma$})}{1+
\exp(\theta_r+\textbf{x}^T\mbox{\boldmath $\gamma$})}, \quad r=1,\cdots, k - 1$$ \\
 $ \iff \displaystyle \log\left\{\frac{P(Y=r|\textbf{x})}{P(Y>r|\textbf{x})}\right\}=
\theta_r+\textbf{x}^T\mbox{\boldmath $\gamma$}.$

\ \

\begin{wideitemize}
\item $\frac{P(Y=r|\textbf{x})}{P(Y>r|\textbf{x})}$ is known as a \textcolor{red}{continuation ratio} \cite{tutz2012}. 

\item Hence, for sequential logit models, the regression coefficients $\boldsymbol \gamma$ have interpretations with respect to this quantity. 
\item (Cf. In the cumulative logistic model, coefficients have interpretations wrt the cumulative odds $\frac{P(Y \leq r | {\bf x})}{P(Y>r|\textbf{x})}$)
\end{wideitemize}
\end{enumerate}
%\end{itemize}
\end{frame}


\begin{frame}
\begin{enumerate}
\setcounter{enumi}{1}
\item
The  exponential CDF \[F(x)=1-e^{-x}\] gives the \textbf{exponential sequential model}:
\begin{eqnarray*}
& P(Y=r|Y\geq r,\textbf{x})=1-e^{-(\theta_r+\textbf{x}^T\mbox{\boldmath $\gamma$})} \\
\iff & -\log\left\{\frac{P(Y>r|\textbf{x})}{P(Y\geq r|\textbf{x})}\right\}=
\theta_r+\textbf{x}^T\mbox{\boldmath $\gamma$}, \quad r=1,\cdots,q.
\end{eqnarray*}
\item
The extreme minimal-value CDF 
\[
F(x)= 1-\exp(-e^x),
\] gives
\begin{equation}
P(Y=r|Y\geq r, \textbf{x})=1-\exp\left(-e^{\theta_r+\textbf{x}^T\mbox{\boldmath $\gamma$}}\right), 
\quad r=1,\cdots, k
\label{eq-3.3.18}
\end{equation}
which is equivalent to
\begin{equation}
\log\left[-\log\left\{\frac{P(Y>r|\textbf{x})}{P(Y\geq r|\textbf{x})}\right\}\right]=
\theta_r+\textbf{x}^T\mbox{\boldmath $\gamma$}, \quad r=1,\cdots, k.
\label{eq-3.3.19}
\end{equation}

\end{enumerate}
\end{frame}
\begin{frame}{\S3.3.4 Extreme min. value CDF $\Rightarrow$ grouped Cox model}
\begin{itemize}
\item
%(\ref{eq-3.3.19}) is equivalent to the cumulative model with 
%$F(x)=1-\exp(-e^x)$. This means 

(\ref{eq-3.3.19}) is a reparametrisation  of the grouped Cox model (a cumulative model). Let $-\infty \equiv \tilde{\theta}_0 < \cdots < \tilde{\theta}_{k-1}$ be such that
$$\theta_r=\log\{e^{\tilde{\theta}_r}-e^{\tilde{\theta}_{r-1}}\}, \quad r=1,\cdots, k-1.$$

\item Note $\theta_1=\tilde{\theta}_1$ and $\displaystyle \tilde{\theta}_r=\log\left(\sum_{i=1}^r e^{\theta_i}\right)$.

\item 
Substituting  $\theta_r=\log\{e^{\tilde{\theta}_r}-e^{\tilde{\theta}_{r-1}}\}$ into (\ref{eq-3.3.19}),
one obtains {\footnotesize
\begin{eqnarray*}
-\log P(Y>r|\textbf{x})+ \log P(Y> r-1|\textbf{x})=e^{\tilde{\theta}_r+\textbf{x}^T\mbox{\boldmath $\gamma$}}
-e^{\tilde{\theta}_{r-1}+\textbf{x}^T\mbox{\boldmath $\gamma$}}, \quad r=1,\cdots,q \\
 \Rightarrow \quad -\log P(Y>r|\textbf{x})=e^{\tilde{\theta}_r+\textbf{x}^T\mbox{\boldmath $\gamma$}},
\quad r=1,\cdots,q \hspace{80pt} \\
 \Rightarrow \quad P(Y\leq r|\textbf{x})=1-\exp\left(-e^{\tilde{\theta}_r+\textbf{x}^T\mbox{\boldmath $\gamma$}}\right),
\quad \mbox{\normalsize  the \textbf{proportional hazard model}.}
\end{eqnarray*}}
\end{itemize}
\end{frame}

\begin{frame}{\S3.3.4 Extended (Generalized) sequential models}

\begin{wideitemize}

\item 

If we assume
\[
\theta_r=\beta_{r0}+{\bf w}^T\boldbeta_r,
\]
 where 
${\boldsymbol \beta}_r$ is a category-specific vector
parameter, one gets the \textit{extended sequential models}:
\begin{equation*}
P(Y=r|Y\geq r,\textbf{x}, \textbf{w})=F(\beta_{r0}+\textbf{w}^T\boldbeta_r+ 
\textbf{x}^T\mbox{\boldmath $\gamma$})
\label{eq-3.3.20}
\end{equation*}
\item $\textbf{w}$ and $\textbf{x}$: respectively called \textcolor{red}{threshold} and \textcolor{red}{shift} variables, as in the extended cumulative models. 


\end{wideitemize}
\end{frame}

\begin{frame}{\large \S3.3.4 Link/response functions and design matrices in sequential models}
\begin{itemize}

\item
The link function $\textbf{g}=(g_1,\cdots, g_q)^T$ is given by
$$g_r(\pi_1,\cdots,\pi_q)=F^{-1}\left(\frac{\pi_r}{1-\pi_1-\cdots -\pi_{r-1}}\right), \quad r=1,\cdots,q$$

\item The response function $\textbf{h}=(h_1,\cdots, h_q)^T$ is given by
$$h_r(\eta_1,\cdots, \eta_q)=F(\eta_r)\prod_{i=1}^{r-1}(1-F(\eta_i)), \quad r=1,\cdots,q.$$
\item
E.g.: For the sequential logit model, with $r=1,\cdots,q$ {\footnotesize
$$g_r(\pi_1,\cdots,\pi_q)=\log\frac{\pi_r}{1\!-\!\pi_1\!-\!\cdots \!-\!\pi_{r}}\quad \mbox{and}\quad
h_r(\eta_1,\cdots, \eta_q)=e^{\eta_r}\prod_{i=1}^{r}(1+e^{\eta_i})^{-1}.$$}

\item
Design matrices structures: No difference between sequential and cumulative models.
Thus the design matrices from \S3.3.3 apply.
\end{itemize}
\end{frame}



\begin{frame}
\begin{itemize}





\item 
The \textcolor{red}{ \texttt{sratio()}} function in the \texttt{R} package VGAM  allows us to perform ordinal regression with sequential models. 

\ \
\item
\cite[Sec. 9.3.3]{tutz2012} compares  cumulative and sequential models in more detail. 
\end{itemize}
\end{frame}
%
%\subsection{\S3.3.5 Strict stochastic ordering}\label{sec:3.3.5}
%\begin{frame}{\S3.3.5 Strict stochastic ordering}
%
%
%
%I leave it for you to read on your own. 
%%\begin{itemize}
%%\item
%%\textbf{Strict stochastic ordering} is a concept generalizing the \textit{proportional odds} for
%%simple cumulative model (\ref{eq-3.3.3})
%%$$\Delta_c(\textbf{x}_1,\textbf{x}_2)=F^{-1}\left(P(Y\leq r |\textbf{x}_1)\right)-
%%F^{-1}\left(P(Y\leq r |\textbf{x}_2)\right).$$
%%If $\Delta_c(\textbf{x}_1,\textbf{x}_2)=\mbox{\boldmath $\gamma$}^T(\textbf{x}_1-\textbf{x}_2)$,
%%we say strict stochastic ordering holds.
%%\item
%%Then for general cumulative model (\ref{eq-3.3.11}) where $\textbf{w}=\textbf{x}$ is set and
%%$\textbf{x}^T\mbox{\boldmath $\gamma$}$ is omitted, one can test strict stochastic ordering by
%%testing $\mbox{\boldmath $\beta$}_1=\mbox{\boldmath $\beta$}_2=\cdots=\mbox{\boldmath $\beta$}_q$.
%%\item
%%The concept can be generalized to the sequential model (\ref{eq-3.3.17}) as well. Let
%%$$\Delta_s(\textbf{x}_1,\textbf{x}_2)=F^{-1}\left(P(Y=r|Y\geq r, \textbf{x}_1)\right)-
%%F^{-1}\left(P(Y=r|Y\geq r, \textbf{x}_2)\right).$$
%%If $\Delta_s(\textbf{x}_1,\textbf{x}_2)=\mbox{\boldmath $\gamma$}^T(\textbf{x}_1-\textbf{x}_2)$,
%%we say strict stochastic ordering holds for the sequential model.
%%\end{itemize}
%\end{frame}
%
%\subsection{\S3.3.6 Two-step models}\label{sec:3.3.6}
%\begin{frame}{\S3.3.6 Two-step models }
%%\begin{itemize}
%%\item
%
%Involves two steps:
%\begin{itemize}
%\item Divide the ordered categories into {\bf sets} of categories with very
%homogeneous responses in the first step. 
%\item Model the categories {\bf within} each set separately
%in the second step. 
%\end{itemize}
%%
%%
%%\end{itemize}
%\end{frame}
%
%
%\begin{frame} {Example 3.9. Patients with Rheumatoid arthritis (Mehta, Patel \& Tsiatis (1984) )}
%
%\begin{itemize}
%%\item Patients
%%with acute rheumatoid arthritis 
%\item New agent compared with an active control (the ``X" in the problem )
%\item Each patient  evaluated on a 5-point assessment scale.  (the ``Y" in the problem )
%\item Global assessment may be subdivided into``improvement", ``no change"
%and ``worse" first. 
%\item Then ``improvement" split up into ``much improved" and ``improved" , and ``worse" category  into ``worse" and ``much worse."
%
%
%\includegraphics{Arthritis}
%%\begin{tabular}{|lccccc|} \hline
%%& \multicolumn{5}{c|}{Global assessment} \\
%%Drug & $\begin{array}{c} \mbox{Much} \\ \mbox{improved}\end{array}$ & Improved & 
%% $\begin{array}{c} \mbox{No} \\ \mbox{change}\end{array}$ & Worse &
%% $\begin{array}{c} \mbox{Much} \\ \mbox{worse}\end{array}$\\ \hline
%%New agent & 24 & 37 & 21 & 19 & 6 \\
%%Active control & 11 & 51 & 22 & 21 & 7 \\ \hline
%%\end{tabular}
%
%%
%%% \vspace*{-8pt}
%%\begin{table} [htbp]
%%\caption{Clinical trial of a new agent and an active control} \label{tab3.8}
%%\begin{center}
%%\begin{tabular}{|lccccc|} \hline
%%& \multicolumn{5}{c|}{Global assessment} \\
%%Drug & $\begin{array}{c} \mbox{Much} \\ \mbox{improved}\end{array}$ & Improved & 
%% $\begin{array}{c} \mbox{No} \\ \mbox{change}\end{array}$ & Worse &
%% $\begin{array}{c} \mbox{Much} \\ \mbox{worse}\end{array}$\\ \hline
%%New agent & 24 & 37 & 21 & 19 & 6 \\
%%Active control & 11 & 51 & 22 & 21 & 7 \\ \hline
%%\end{tabular}
%%\end{center}
%%\end{table}
%
%\end{itemize}
%\end{frame}
%
%
%
%
%\begin{frame}{\S3.3.6 Two-step models formulation}
%\begin{itemize}
%\item
%The categories $1, \cdots, k$ are subdivided into $t$ basic sets $S_1, \cdots, S_t$,
%where $S_j=\{m_{j-1}+1, \cdots, m_j\}$, $m_0=0$ and $m_t=k$. \vspace{8pt}
%
%In \underline{Step 1}: $Y\in S_j\;\; \Leftrightarrow \;\; \theta_{j-1}<U_0\leq \theta_j,\;\;
%j=1,\cdots, t.$
%
%\hspace{50pt} Assume $U_0=-\textbf{x}^T\mbox{\boldmath $\gamma$}_0+\varepsilon$.
%\vspace{6pt}
%
%In \underline{Step 2}: $Y=R \mid Y\in S_j \;\; \Leftrightarrow \;\; \theta_{j, r-1}<U_j\leq \theta_{jr}.$
%
%\hspace{50pt} Assume $U_j=-\textbf{x}^T\mbox{\boldmath $\gamma$}_j+\varepsilon_j$.
%Also $\varepsilon, \varepsilon_j$'s are iid with cdf $F$.
%\item
%Then the two-step model becomes
%\begin{equation}
%P(Y\in T_j|\textbf{x})=F(\theta_j+\textbf{x}^T\mbox{\boldmath $\gamma$}_0), \quad
%P(Y\leq r |Y\in S_j, \textbf{x})=F(\theta_{jr}+\textbf{x}^T\mbox{\boldmath $\gamma$}_j)
%\label{eq-3.3.21}
%\end{equation}
%where $T_j=S_1\cup\cdots\cup S_j, \quad \theta_1<\cdots <\theta_{t-1}, \quad \theta_t=\infty$;
%$\theta_{j, m_{j-1}+1}<\cdots <\theta_{j, m_{j-1}}, \quad \theta_{j, m_j}=\infty, \quad j=1,\cdots,t$.
%
%
%\end{itemize}
%\end{frame}
%
%\begin{frame}
%
%When cumulative mechanisms are used in both steps, (\ref{eq-3.3.21}) is called a \textbf{two-step 
%cumulative model}. \textbf{Two-step sequential model} can be similarly defined.
%
%\ \ 
%
%This is still in the framework of GLMs. I leave it to you to read about the forms of the link function, design matrix etc. 
%\end{frame}
%
%\subsection{\S3.3.7 Alternative approaches}\label{sec:3.3.7}
%\begin{frame}{\S3.3.7 Alternative approaches (Not covered in detail)}
%\begin{enumerate}
%\item
%{\bf Stereotype regression}, Anderson (1984):
%$$P(Y=r|\textbf{x})=\frac{e^{\beta_{r0}-\phi_r\mbox{\scriptsize \boldmath 
%$\beta$}^T\textbf{x}}}{1+\sum_{i=1}^q e^{\beta_{i0}-\phi_i\mbox{\scriptsize \boldmath $\beta$}^T\textbf{x}}},\quad
%r=1,\cdots, q;$$
%where $1=\phi_1>\cdots >\phi_k=0$.
%\item
%Williams and Grizzle (1972).
%$\displaystyle \sum_{r=1}^k s_rP(Y=r|\textbf{x})=\textbf{x}^T\mbox{\boldmath $\beta$}$
%
%where $s_1,\cdots, s_k$ are some given scores for the categories.
%\item
%\textbf{Adjacent categories logit} (Agresti 1984):
%\begin{eqnarray*}
%\lefteqn{\log \frac{P(Y=r|\textbf{x})}{P(Y=r-1|\textbf{x})}=\textbf{x}^T\mbox{\boldmath $\beta$}_r} & \\
%&& \Leftrightarrow \quad P(Y=r|Y\in\{r, r+1\}, \textbf{x})=F(\textbf{x}^T\mbox{\boldmath $\beta$}_r)
%\end{eqnarray*}
%where $F$ is the logistic cdf.
%\end{enumerate}
%\end{frame}
%
\section{\S3.4 Statistical inference for MGLM}\label{sec:3.4}
\begin{frame}{\S3.4 Statistical inference for MGLM overview}
\begin{enumerate}
\item
Maximum likelihood estimation
\item
Testing of linear hypotheses
\item
Goodness of fit (model adequacy) statistics
\item
Power-divergence statistics
\end{enumerate}
\end{frame}

\subsection{\S3.4.1 Maximum likelihood estimation in MGLM}\label{sec:3.4.1}
%\begin{frame}{\S3.4.1 Maximum likelihood estimation in MGLM (1)}
%\begin{itemize}
%\item
%An extension of the exponential family is the \textbf{multivariate exponential family}, having the pdf
%$$f(\textbf{y}_i | \mbox{\boldmath $\theta$}_i, \phi, \omega_i)=\exp\left\{\frac{\textbf{y}_i^T
% \mbox{\boldmath $\theta$}_i-b(\mbox{\boldmath $\theta$}_i)}{\phi}\cdot \omega_i + 
%c(\textbf{y}_i,\phi, \omega_i)\right\}$$
%\item
%For $q\times 1$ vector observations $\textbf{y}_1,\cdots, \textbf{y}_n$ which are suppose to be independent
%of each other, let $\mbox{\boldmath $\mu$}_i=E(\textbf{y}_i|\textbf{x}_i)$, $i=1,\cdots,n$, with $\textbf{x}_i$ being the 
%covariate vector.
%\item
%The \textbf{structural assumption} of MGLM is that $\mbox{\boldmath $\mu$}_i$'s depend on the linear
%predictor  $\mbox{\boldmath $\eta$}_i=Z_i\mbox{\boldmath $\beta$}$ in the form
%$$\mbox{\boldmath $\mu$}_i=\textbf{h}(\mbox{\boldmath $\eta$}_i)=
%\textbf{h}(Z_i\mbox{\boldmath $\beta$})=\left[h_1(Z_i\mbox{\boldmath $\beta$}),
%\cdots,h_1(Z_i\mbox{\boldmath $\beta$})\right]^T, \quad i=1,\cdots,n$$
%where
%\begin{itemize}
%\item
%the response function $\textbf{h}: S\rightarrow M$ is defined on $S\subset \mathbb{R}^q$, taking
%values in $M\subset\mathbb{R}^q$;
%\item
%$Z_i$ is a $q\times p$ design matrix for unit $i$;
%\item
%$\mbox{\boldmath $\beta$}=(\beta_1,\cdots,\beta_p)^T$ is an unknown parameter vector
%from $B\subset\mathbb{R}^q$.
%\end{itemize}
%\end{itemize}
%\end{frame}

\begin{frame}{Some convention for derivatives}
 In what follows, for a function ${\bf f} :  \mathbb{R}^q\longrightarrow  \mathbb{R}^q$, where 
 \[
 {\bf f}: {\bf t}= (t_1, \dots, t_q)  \mapsto (f_1({\bf t}), \cdots, f_q({\bf t}))^T,
 \]
we let 
\begin{itemize}
\item
\[
\frac{\partial {\bf f}}{\partial {\bf t}} =  
\begin{bmatrix}
\frac{\partial f_1}{\partial t_1} & \cdots & \frac{\partial f_1}{\partial t_q}\\
\vdots & \vdots & \vdots\\
\frac{\partial f_q}{\partial t_1} & \cdots &  \frac{\partial f_q}{\partial t_q}
\end{bmatrix}
\]

\item
\[
\frac{\partial {\bf f}^T}{\partial {\bf t}} =  
\begin{bmatrix}
\frac{\partial f_1}{\partial t_1} & \cdots & \frac{\partial f_q}{\partial t_1}\\
\vdots & \vdots & \vdots\\
\frac{\partial f_1}{\partial t_q} & \cdots &  \frac{\partial f_q}{\partial t_q}
\end{bmatrix}
\]
\end{itemize}

\end{frame}

\begin{frame}{\S3.4.1 Maximum likelihood estimation in MGLM}
\begin{itemize}
\item
The MLE of $\mbox{\boldmath $\beta$}$ may be derived in analogy to the one-dimensional case.
% an item
\item
The log-likelihood ``kernel" for observation $\textbf{y}_i$ is
\[
\ell_i(\boldmu_i)=\frac{{\bf y}_i^T\boldtheta_i
-b(\boldtheta_i)}{\phi}\cdot \omega_i, \quad \mbox{where $\boldtheta_i
=\boldtheta(\boldmu_i)$}
\]

%another item
\item
Using  $\boldmu_i=\textbf{h}(\boldeta_i)=
\textbf{h}(Z_i\boldbeta)$, by the chain rule, the score function is  
\[
\textbf{s}(\boldbeta)\equiv\frac{\partial\ell}{\partial \boldbeta}=
\sum_{i=1}^n \textbf{s}_i(\boldbeta), \text{ for }
\] 
\begin{enumerate}
\item $\textbf{s}_i(\boldbeta)\equiv Z_i^TD_i(\boldbeta) \underbrace{\Sigma_i^{-1}}_{\frac{\partial \boldtheta}{ \partial \boldmu}}(\boldbeta)
[\textbf{y}_i-  \underbrace{\boldmu_i(\boldbeta)}_{ \frac{\partial b}{\partial \boldtheta}}]$
\item $D_i(\boldbeta)\equiv
\frac{\partial\textbf{h}^T(\boldeta_i)}{\partial \boldeta_i}$,  the
(square) derivative {\bf matrix} 
 evaluated at $\boldeta_i
=Z_i\boldbeta
$, 
\item  $\Sigma_i (\boldbeta)\equiv \text{cov}(\textbf{y}_i) = \frac{\phi}{\omega_i} \frac{\partial^2 b(\boldtheta)}{\partial \boldtheta^2}$.
\end{enumerate}
\end{itemize}
\end{frame}

\begin{frame}{More on the score function}
\begin{wideitemize}

\item
An alternative form for each $\textbf{s}_i(\mbox{\boldmath $\beta$})$ is
\[
\textbf{s}_i(\boldbeta)=Z_i^T \cdot W_i(\boldbeta)\cdot \frac{\partial
\textbf{g}^T(\boldmu_i)}{\partial \boldmu_i}
\cdot [\textbf{y}_i-\boldmu_i(\boldbeta)],
\]
where $W_i$ is a weight matrix defined as
\begin{equation*}
W_i(\boldbeta)\equiv D_i(\boldbeta)\Sigma_i^{-1}(\boldbeta)
D^T_i(\boldbeta)=\left\{\frac{\partial\textbf{g}(\boldmu_i)}{\partial 
\boldmu_i} \Sigma_i(\boldbeta)\frac{\partial
\textbf{g}^T(\boldmu_i)}{\partial \boldmu_i}\right\}^{-1}.
\end{equation*}

\item (Cf.  $w_i (\boldbeta) = D_i^2\sigma_i^{-2} $ in the univariate case.)


\item $W_i$ approximately equals $(\mbox{cov}(\textbf{g}(\textbf{y}_i))^{-1}$, by a linearization argument.

\item $\frac{\partial\textbf{g}^T(\boldmu_i)}{\partial 
\boldmu_i}$ is the inverse of $D_i(\mbox{\boldmath $\beta$}) = \frac{\partial\textbf{h}^T(\boldeta_i)}{\partial \boldeta_i}$, by inverse function theorem.


\end{wideitemize}
\end{frame}

\begin{frame} {Fisher informaton}

The (expected) Fisher information matrix is $$F(\mbox{\boldmath $\beta$})=
\text{cov}\Big(\textbf{s}(\mbox{\boldmath $\beta$})\Big)=\sum_{i=1}^nZ_i^TW_i(\mbox{\boldmath $\beta$})Z_i.$$
\end{frame}

\begin{frame}{Score and Fisher information in matrix form}
\begin{wideitemize}
\item
In pure matrix form, the above may be rewritten as
\begin{eqnarray*}
\textbf{s}(\mbox{\boldmath $\beta$}) &=& Z^TD(\mbox{\boldmath $\beta$})\Sigma^{-1}(\mbox{\boldmath $\beta$})
[\textbf{y}-\mbox{\boldmath $\mu$}(\mbox{\boldmath $\beta$})] \\
F(\mbox{\boldmath $\beta$}) &=& Z^TW(\mbox{\boldmath $\beta$})Z.
\end{eqnarray*}
where 
\begin{itemize}
\item 
$\textbf{y}=(\textbf{y}_1^T,\cdots, \textbf{y}_n^T)^T$,
$\mbox{\boldmath $\mu$}(\mbox{\boldmath $\beta$})=(\mbox{\boldmath $\mu$}_1(\mbox{\boldmath $\beta$})^T,
\cdots,\mbox{\boldmath $\mu$}_n(\mbox{\boldmath $\beta$})^T)^T$
\item Block diagonal matrices:
\begin{multline*}
\Sigma(\mbox{\boldmath $\beta$})=\mbox{diag}\left(\Sigma_i(\mbox{\boldmath $\beta$})\right),\\
W(\mbox{\boldmath $\beta$})=\mbox{diag}\left(W_i(\mbox{\boldmath $\beta$})\right) \text{ and }\\
D(\mbox{\boldmath $\beta$})=\mbox{diag}\left(D_i(\mbox{\boldmath $\beta$})\right).
\end{multline*}

\item  $\displaystyle Z=\left[Z_1^T, \cdots, Z_n^T\right]^T$
\end{itemize}
\item
For grouped data, $\textbf{y}_i$ is the mean vector over $n_i$ observations and
$\Sigma_i(\mbox{\boldmath $\beta$})$ is replaced by $n_i^{-1}\Sigma_i(\mbox{\boldmath $\beta$})$.

\end{wideitemize}
\end{frame}

\begin{frame} {Standard properties of MLE}
\begin{wideitemize}
\item
Under regularity conditions, the MLE is asymptotically normal:
$$\; \hat{\mbox{\boldmath $\beta$}}\stackrel{a}{\sim}
N(\mbox{\boldmath $\beta$}, F^{-1}(\hat{\mbox{\boldmath $\beta$}})).$$


\item More precisely, for $F_\infty (\boldbeta) =  \lim_n \frac{1}{n} F(\boldbeta) $,
\[
\sqrt{n} (\hat{\boldsymbol \beta} - {\boldsymbol \beta}) \longrightarrow_d N(0, F_\infty^{-1} ({\boldsymbol \beta}))
\]
in distribution. 
\item See Appendix A2 in F \& T for details.
\end{wideitemize}

\end{frame}
%
%\begin{frame} {Multivariate Delta Method (Digression)}
%
%\begin{itemize}
%%\item A general  technique to derive standard errors (Not just for  GLM)
%
%\item For the moment, let $\phi$, $\theta$ and $\Sigma$ be quantities that have nothing to do with the parameters of a GLM.
%
%\item  Suppose  
%
%\begin{itemize}
%\item $T_n = (T_{n1}, \dots, T_{np})^T$ is a random vector such that 
%\[
%\sqrt{n}(T_n -  \theta ) \longrightarrow N(0, \Sigma)
%\]
%
%
%\item A function $g (t_1, \dots, t_p)$ of $t = (t_1,\dots, t_p)^T$ has a \emph{non-zero} derivative $\phi = (\phi_1, \dots, \phi_p)^T$ at $\theta$, i.e.
%\[
%\phi_i = \frac{\partial g}{\partial t_i} |_{t = \theta}
%\]
%\end{itemize}
%\item Then,
%\[
%\sqrt{n}(g(T_n) - g(\theta)) \rightarrow_d N(0,   \phi^T \Sigma \phi)
%\]
%
%
%
%
%\end{itemize}
%
%
%\end{frame}
%
%
%
%
%% new frame
%\begin{frame} {Multivariate Delta Method (Digression)}
%\begin{wideitemize}
%
%
%\item In practice, $\phi = \phi(\theta)$ and/or $\Sigma =\Sigma(\theta)$ may depend on the unknown parameter $\theta$ itself. One may need to substitute a consistent estimate for $\theta$ (such has $T_n$ itself) into $\phi^T \Sigma \phi$ to form the covariance estimate. 
%\item Implication in the GLM context: If $g(\cdot)$ a differentiable function on ${\boldsymbol \beta}$, then
%\[
%\sqrt{n} (g(\hat{\boldsymbol \beta}) -g( {\boldsymbol \beta})) \longrightarrow_d N\Big(0,   \partial g^T(\hat{\boldsymbol \beta}) F_\infty^{-1} ({\boldsymbol \beta}) g(\hat{\boldsymbol \beta}) \Big).
%\]
%In other words, 
%\[
%g(\hat{\boldsymbol \beta})\stackrel{a}{\sim}  N(g({\boldsymbol \beta}), \partial g^T(\hat{\boldsymbol \beta}) F^{-1} (\hat{ \boldsymbol \beta}) \partial g(\hat{\boldsymbol \beta}) ).
%\] 
%
%\end{wideitemize}
%\end{frame}

\begin{frame}{\S3.4.1 Fisher scoring in MGLM}

Same as in the univariate case, except using the multivariate versions of vectors and matrices.
\begin{wideitemize}
\item
The working or pseudo observation vector is
\[
\tilde{\bf y}(\boldbeta) \equiv (\tilde{\bf y}_1(\boldbeta)^T,
\cdots, \tilde{\bf y}_n(\boldbeta)^T)^T,
\]
where
\[
 \tilde{\bf y}_i(\boldbeta)= Z_i\boldbeta+
\left(D_i^{-1}(\boldbeta)\right)[{\bf y}_i-\boldmu_i(\boldbeta)]
\]
is a Taylor approximation for $\textbf{g}(\mbox{\boldmath $y$}_i(\mbox{\boldmath $\beta$}))$ around $\boldmu_i(\boldbeta) = Z_i \boldbeta$.
\item
The iterated reweighted least square (IRLS) estimate is
$$\hat{\mbox{\boldmath $\beta$}}^{(k+1)}=\left(Z^TW(\hat{\mbox{\boldmath $\beta$}}^{(k)})Z\right)^{-1}
Z^TW(\hat{\mbox{\boldmath $\beta$}}^{(k)})\tilde{\textbf{y}}(\mbox{\boldmath $\beta$}^{(k)})$$
which is equivalent to the Fisher scoring iteration
$$\hat{\mbox{\boldmath $\beta$}}^{(k+1)}=\hat{\mbox{\boldmath $\beta$}}^{(k)}+
\left(Z^TW(\hat{\mbox{\boldmath $\beta$}}^{(k)})Z\right)^{-1}\textbf{s}(\mbox{\boldmath $\beta$}^{(k)}).$$
\end{wideitemize}

\end{frame}



\section{\S3.4 Statistical Inference} \label{sec:3.4}

\subsection{\S3.4.2 Testing and goodness-of-fit in MGLM}\label{sec:3.4.2}
\begin{frame}{\S3.4.2 Testing linear hypotheses in MGLM}
The linear hypotheses
\[
H_0: C\boldbeta=\boldxi\quad \mbox{vs.}\quad
H_1: C \boldbeta\neq\boldxi
\]
can apparently be tested by the trilogy of classical tests. Replace the score and
Fisher information in the univariate GLM  by their multivariate versions.
\end{frame}

\begin{frame} {Goodness of fit statistics for MGLM}

Assume: data are grouped as far as possible

\begin{itemize}
\item Pearson statistic
\item Deviance statistic
\end{itemize}

\end{frame}

\begin{frame}{Pearson statistic in MGLM}

\begin{itemize}
\item
\textbf{Pearson Statistic} for MGLM has the generic form:
\[
\chi^2 =  \sum_{i=1}^g (\textbf{y}_i-\hat{\mbox{\boldmath $\mu$}}_i)^T
\Sigma_i^{-1}(\hat{\boldsymbol \beta})(\textbf{y}_i-\hat{\mbox{\boldmath $\mu$}}_i).
\]

\item In this definition, $ \hat{\Sigma}_i  := \Sigma_i(\hat{\boldsymbol \beta})$ is assumed to be invertible.

% ${\bf y}_i$ is a vector of  $k$ dependent relative frequencies, so 

\end{itemize}
\end{frame}
% new frame
\begin{frame}{Pearson statistic for multinomial vector}
\begin{wideitemize}
\item 
 If
${\bf y}_i$  is a \underline{scaled} multinomial vector of length $q = k-1$, where $k$ is the no.  of categories, i.e.
\[
n_i\textbf{y}_i\sim
M(n_i, \boldpi_i) \text{ for } \boldpi_i = (\pi_{i1}, \dots, \pi_{iq})^T,
\]
 the Pearson statistic will have the familiar form (cf. MAST20005)
\begin{equation} \label{Pearson_multinom_form}
\chi^2=\sum_{i=1}^g \chi_p^2(\textbf{y}_i, \hat{\mbox{\boldmath $\pi$}}_i)=
\sum_{i=1}^g n_i\sum_{j=1}^k\frac{(y_{ij}-\hat{\pi}_{ij})^2}{\hat{\pi}_{ij}},
\end{equation}
where 
\begin{itemize}
\item $\hat{\pi}_{ij} = \pi_{ij}(\hat{{\boldsymbol \beta}})$ for $j = 1, \dots, q$
\item  $y_{ik}=1-y_{i1}-\cdots-y_{iq}$
\item   $\hat{\pi}_{ik}=1-\hat{\pi}_{i1}-\cdots-\hat{\pi}_{iq}$
\end{itemize}



\item To derive \eqref{Pearson_multinom_form} from the generic form of Pearson statistic, one needs to invert
 $
 \hat{\Sigma}_i 
$
for the scaled multinomial vector $\textbf{y}_i$.  

\end{wideitemize}
\end{frame}

% new frame

\begin{frame}{Inverting $\Sigma_i$}
\begin{wideitemize}
 \item Recall for $\boldpi_i  = {\bf h}({\boldsymbol \eta}_i) = {\bf h}(Z_i{\boldsymbol \beta})$, 
\[
\Sigma_i =\frac{\left(\mbox{diag}(\mbox{\boldmath $\pi$}_i)-\mbox{\boldmath $\pi$}_i
\mbox{\boldmath $\pi$}_i^T\right)}{n_i},
\]
which is the \underline{sum of a diagonal and  rank-one matrix}.

\item The Woodbury matrix identity is useful: For any conformable  matrices $A, U, C$  and $V$,
\[
(A + UCV)^{-1} = A^{-1} - A^{-1}  U(C^{-1} + V A^{-1}U)^{-1} V A^{-1}
\]



\end{wideitemize}
\end{frame}


% new frame

\begin{frame}{Pearson statistic for multinomial vector, cont'd}
\begin{wideitemize}
\item  Hence, taking $A = \text{diag}(\boldpi_i)$, $C = -1$, $U = \boldpi_i$, $V =  \boldpi_i^T$,
\begin{align*}
&\left(\mbox{diag}(\mbox{\boldmath $\pi$}_i)-\mbox{\boldmath $\pi$}_i
\mbox{\boldmath $\pi$}_i^T\right)^{-1} \\
&= \text{diag}(\boldpi_i)^{-1} -  \text{diag}(\boldpi_i)^{-1} \boldpi_i (  \boldpi_i^T \text{diag}(\boldpi_i)^{-1}  \boldpi_i - 1 )^{-1}  \boldpi_i^T \text{diag}(\boldpi_i)^{-1}\\
&= \text{diag}(\boldpi_i)^{-1} +\pi_{ik}^{-1} {\bf 1}_q {\bf 1}_q^T.
\end{align*}


\item So 
\begin{align*}
& (\textbf{y}_i- \boldpi_i)^T
\Sigma_i^{-1}  (\textbf{y}_i- \boldpi_i)\\
&= n_i \Big[(\textbf{y}_i- \boldpi_i)^T \text{diag}(\boldpi_i)^{-1}(\textbf{y}_i- \boldpi_i) + \pi_{ik}^{-1} (\textbf{y}_i- \boldpi_i)^T {\bf 1}_q {\bf 1}_q^T  (\textbf{y}_i- \boldpi_i) \Big] \\
&= n_i \bigg[\sum_{j=1}^q\frac{(y_{ij}-{\pi}_{ij})^2}{{\pi}_{ij}} + \pi_{ik}^{-1} (\pi_{ik} - y_{ik} )^2\bigg].
\end{align*}

\item Replacing $\pi_1, \dots, \pi_k$ with $\hat{\pi}_1, \dots, \hat{\pi}_k$ gives \eqref{Pearson_multinom_form}.
\end{wideitemize}


\end{frame}



% new frame
%\begin{frame} {Pearson statistic for multinomial vector, cont'd}
%\begin{wideitemize}
%
%
%\item 
%
%
%The simple form \eqref{Pearson_multinom_form} can be more easily derived by constructing the Pearson statistic with the scaled multinomial vector of \underline{full length $k$},
%\[
%\mathring{\bf y}_i \equiv(y_{i1}, \dots, y_{ik})^T,
%\]
% which has a \underline{non-invertible} $k$-by-$k$ covariance matrix, call it $\mathring{\Sigma}_i$.
%
%\item $\Sigma_i$ is the upper-left $q$-by-$q$ submatrix of $\mathring{\Sigma}_i$.
%\item  In this case, one need to form the Pearson statistic with a  \emph{generalized inverse} (g-inverse) $\mathring{\Sigma}_i^-$, a matrix such that 
%\[
%\mathring{\Sigma}_i \mathring{\Sigma}_i^- \mathring{\Sigma}_i= \mathring{\Sigma}_i.
%\] 
%Note that, a g-inverse is not unique in general.
%
% 
%  
%
%
%
%\end{wideitemize}
%
%\end{frame}
%
%\begin{frame}
%
%\begin{wideitemize}
%
%\item In particular, if we take $\mathring{\Sigma}_i^- 
% = \begin{bmatrix}
%     \Sigma_i^{-1} & {\bf 0}           \\
%     {\bf 0}       ^T & 0          
%     \end{bmatrix}  $; one can easily check that it is indeed a g-inverse for $\mathring{\Sigma}_i$. 
%     \item Moreover,    for $ \hat{\mathring{\boldpi}}_i = (\hat{\pi}_{i1}, \dots, \hat{\pi}_{ik})^T$, 
% \[
% \sum_{i = 1}^g(\mathring{\bf y}_i -  \hat{\mathring{\boldpi}}_i)^T
% \Big[ \begin{smallmatrix}
%     \hat{\Sigma}_i^{-1} & {\bf 0}           \\
%     {\bf 0}       ^T & 0          
%     \end{smallmatrix}\Big]
%  (\mathring{\bf y}_i -  \hat{\mathring{\boldpi}}_i) = \sum_{i = 1}^g({\bf y}_i -  \hat{\boldpi}_i)^T\hat{{\Sigma}}_i^{-1} ({\bf y}_i -  \hat{\boldpi}_i).
% \]
%
% 
% 
% \item
% Unfortunately, we still need to invert ${\Sigma}_i$ to obtain this $\mathring{\Sigma}_i^-$.
%
%
%\end{wideitemize}
%
%
%\end{frame}
%
%%new frame
%
%
%\begin{frame} {Pearson statistic for multinomial vector, cont'd}
%
%\begin{wideitemize}
%\item Alternatively, by taking  the $\mathring{\Sigma}_i^-$  as 
% \[
%\mathring{ \Lambda}_i  = n_i\text{diag}(\pi_{i1}^{-1} , \dots,   \pi_{ik}^{-1}),
% \]
%which can be easily verified to be a g-inverse of $\mathring{\Sigma}_i$, the construction
% \[
% \sum_{i = 1}^g(\mathring{\bf y}_i -  \hat{\mathring{\boldpi}}_i)^T
%\hat{\mathring{ \Lambda}}_i
%  (\mathring{\bf y}_i -  \hat{\mathring{\boldpi}}_i)
%   \]
%  easily gives  \eqref{Pearson_multinom_form}, where $\hat{\mathring{ \Lambda}}_i = n_i\text{diag}(\hat{\pi}_{i1}^{-1} , \dots,   \hat{\pi}_{ik}^{-1})$.
%
%\item Turns out, \underline{ all choices of the g-inverse will lead to \eqref{Pearson_multinom_form}}. 
%
%
%\end{wideitemize}
%
%
%\end{frame}
%
%
%\begin{frame}{Pearson statistic for multinomial vector, cont'd}
%\begin{wideitemize}
%\item By  \cite[Theorem 2.1]{radhakrishna1972generalized}, any g-inverse of $\mathring{\Sigma}_i$ has the generic form
%\[
%\mathring{\Sigma}_i^- + U - \mathring{\Sigma}_i^- \mathring{\Sigma}_i U \mathring{\Sigma}_i \mathring{\Sigma}_i^- , 
%\]
%where $U$ is an arbitrary matrix of matching dimensions. The $\mathring{\Sigma}_i^-$ in display can be any one particular g-inverse of  $\mathring{\Sigma}_i$. 
%
%
%
% \item  One can compute 
%\begin{multline*}
%\sum_{i = 1}^g(\mathring{\bf y}_i -  \hat{\mathring{\boldpi}}_i)^T \left(  U -   
%\widehat{\ringLambda}_i
%     \widehat{\mathring{\Sigma}}_i U  \widehat{\mathring{\Sigma}}_i 
%    \widehat{\ringLambda}_i  \right) (\mathring{\bf y}_i -  \hat{\mathring{\boldpi}}_i) = \\
%  \sum_{i = 1}^g (\mathring{\bf y}_i -  \hat{\mathring{\boldpi}}_i)^T U (\mathring{\bf y}_i -  \hat{\mathring{\boldpi}}_i)   - \sum_{i = 1}^g  \Big(\widehat{\mathring{\Sigma}}_i 
%    \widehat{\ringLambda}_i  (\mathring{\bf y}_i -  \hat{\mathring{\boldpi}}_i) \Big) ^T  U \Big(\widehat{\mathring{\Sigma}}_i 
%    \widehat{\ringLambda}_i  (\mathring{\bf y}_i -  \hat{\mathring{\boldpi}}_i) \Big)\textcolor{red}{ = 0},
% \end{multline*}
%because 
%\[
% \widehat{\mathring{\Sigma}}_i 
%    \widehat{\ringLambda}_i  (\mathring{\bf y}_i -  \hat{\mathring{\boldpi}}_i) = \mathring{\bf y}_i -  \hat{\mathring{\boldpi}}_i,
%    \]
%    as will be shown below:
%
%\end{wideitemize}
%
%\end{frame}
%
%
%\begin{frame} {  $ \widehat{\mathring{\Sigma}}_i 
%    \widehat{\ringLambda}_i  (\mathring{\bf y}_i -  \hat{\mathring{\boldpi}}_i) = \mathring{\bf y}_i -  \hat{\mathring{\boldpi}}_i$}
%
% First 
%
%
%\[
%\mathring{\Sigma}_i = 
%n_i^{-1}
%\left[
%\begin{array}{cccc} \pi_{i1}(1\!-\!\pi_{i1}) & -\pi_{i1}\pi_{i2} & \cdots
%& -\pi_{i1}\pi_{ik} \\
% -\pi_{i2}\pi_{i1} & \pi_{i2}(1\!-\!\pi_{i2}) & \cdots & -\pi_{i2}\pi_{ik} \\
%  \vdots & \vdots & \ddots 
%& \vdots \\ -\pi_{ik}\pi_{i1} & -\pi_{ik}\pi_{i2} & \cdots & \pi_{ik}(1\!-\!\pi_{ik})
%\end{array}
%\right].
%\]
%So we have
%\[
%\mathring{\Sigma}_i  \ringLambda_i   = 
%\left[
%\begin{array}{cccc} 1\!-\!\pi_{i1} & -\pi_{i1} & \cdots
%& -\pi_{i1} \\
% -\pi_{i2} & 1\!-\!\pi_{i2} & \cdots & - \pi_{i2} \\
%  \vdots & \vdots & \ddots 
%& \vdots \\ -\pi_{ik}& -\pi_{ik}& \cdots & 1\!-\!\pi_{ik}
%\end{array}
%\right].
%\]
%
%
%\end{frame}
%
%
%\begin{frame} {  $ \widehat{\mathring{\Sigma}}_i 
%    \widehat{\ringLambda}_i  (\mathring{\bf y}_i -  \hat{\mathring{\boldpi}}_i) = \mathring{\bf y}_i -  \hat{\mathring{\boldpi}}_i$, cont'd}
%\begin{wideitemize}
%\item
%Continuing, we get that
%\begin{align*}
%\mathring{\Sigma}_i \ringLambda_i   (\mathring{\bf y}_i -  \mathring{\boldpi}_i) &= 
%\left[
%\begin{array}{cccc} 1\!-\!\pi_{i1} & -\pi_{i1} & \cdots
%& -\pi_{i1} \\
% -\pi_{i2} & 1\!-\!\pi_{i2} & \cdots & - \pi_{i2} \\
%  \vdots & \vdots & \ddots 
%& \vdots \\ -\pi_{ik}& -\pi_{ik}& \cdots & 1\!-\!\pi_{ik}
%\end{array}
%\right] 
%\left[
%\begin{array}{c}
%y_{i1} - \pi_{i1}\\
%y_{i2} - \pi_{i2}\\
%\vdots\\
%y_{ik} - \pi_{ik}
%\end{array}
%\right] 
%\\
%&=
%\left[
%\begin{array}{c}
%y_{i1} - \pi_{i1} - \pi_{i1} \sum_{j=1}^k (y_{ij} - \pi_{ij})\\
%y_{i2} - \pi_{i2} - \pi_{i2} \sum_{j=1}^k (y_{ij} - \pi_{ij}) \\
%\vdots\\
%y_{ik} - \pi_{ik} - \pi_{ik} \sum_{j=1}^k (y_{ij} - \pi_{ij})
%\end{array}
%\right]  =  \left[
%\begin{array}{c}
%y_{i1} - \pi_{i1} \\
%y_{i2} - \pi_{i2} \\
%\vdots\\
%y_{ik} - \pi_{ik}
%\end{array}
%\right],
%\end{align*}
%because $\sum_{j=1}^k (y_{ij} - \pi_{ij} )=  \sum_{j=1}^k y_{ij} - \sum_{j=1}^k   \pi_{ij} = 1- 1 = 0$. 
%
%\item Finally,
%replace $\mathring{\Sigma}_i$,   $\ringLambda_i $ and $ \mathring{\boldpi}_i$ with their  estimates by plugging in $\hat{\boldsymbol \beta}$.
%\end{wideitemize}
%\end{frame}

\begin{frame}{Deviance}
\begin{wideitemize}
\item For multinomial data, one can derive the deviance as 

$$D=-2\sum_{i=1}^g\left\{\ell_i(\hat{\mbox{\boldmath $\pi$}}_i)-\ell_i(\textbf{y}_i)\right\},$$
or more explicitly,

$$D=2\sum_{i=1}^g\chi_D^2(\textbf{y}_i, \hat{\mbox{\boldmath $\pi$}}_i)=
2\sum_{i=1}^g n_i\sum_{j=1}^k y_{ij}\log\frac{y_{ij}}{\hat{\pi}_{ij}}.
$$

where $p$ is the number of estimated parameters.

\item Note: Since the dispersion $\phi = 1$ for multinomial data, the deviance in this case is same as the \emph{scaled} deviance, which is twice the log-likelihood difference between the saturated and GLM model. 
%\item
%For sparse data with small $n_i$, one should use alternative asymtotics.
\end{wideitemize}
\end{frame}


\begin{frame} {Chi-square asymptotics for the goodness-of-fit statistics}

\begin{itemize}
\item
Under regularity conditions (fixed $g$, $n_i/n \rightarrow \lambda_i > 0$), the above goodness-of-fit statistics for multinomial will approximately follow
\[
\chi^2, D\stackrel{a}{\sim} \chi^2(g(k-1)-p)
%\quad \mbox{and} \quad
%D\stackrel{a}{\sim}\chi^2(g(k-1)-p)
\]
when the model is true. 
\end{itemize}


\end{frame}


\begin{frame} {Example 3.11 (Caesarian birth study cont'd)}
\begin{itemize}
\item Recall we have a $3 \times 2 \times 2 \times 2$ contingency table for the variables 
\begin{enumerate}
\item response: infection  type (3 levels: type 1, type 2, no infection)
\item covariates: NOPLAN, FACTOR, ANTIB (each with 2 levels)
\end{enumerate}
\item The two equations are 
\begin{eqnarray*}
\log\frac{P(\mbox{infection type I})}{P(\mbox{no infection})} 
=\beta_{10}+\beta_{1N}\texttt{NoPlan}+\beta_{1A}\texttt{Antib}+\beta_{1R}\texttt{RiskF} \\
\log\frac{P(\mbox{infection type II})}{P(\mbox{no infection})} 
=\beta_{20}+\beta_{2N}\texttt{NoPlan}+\beta_{2A}\texttt{Antib}+\beta_{2R}\texttt{RiskF}
\end{eqnarray*}

\item F \&T (p.109) reported a deviance of 11.830 for this model.


\end{itemize}
\end{frame}


\begin{frame}
 
 \begin{wideitemize}
 
\item F \& T also tests $H_0 : \beta_{1N} = \beta_{2N} \text{ and } \beta_{1F} = \beta_{2F} $ against the full model above, with likelihood ratio test (limiting $\chi^2$ with d.f $2$). 
 
 \item This requires us to perform MLE under the constraints $ \beta_{1N} = \beta_{2N} \text{ and } \beta_{1F} = \beta_{2F} $, which 
 can be done with \texttt{vglm} in the package \texttt{VGAM}. The book reported a deviance of 12.162 for this restricted model.
 
 
 \item Likelihood ratio statistics can be then computed as 
 \[
 12.162 - 11.830 = 0.332.
 \]
 (F \&T p.109 misprinted this number to be 0.8467)

 \item p-value is 0.847, not significant.
 \end{wideitemize}

\end{frame}


\subsection{\S3.4.3 Power-divergence family}\label{sec:3.4.3}
%\begin{frame}
%
%
%\end{frame}


\begin{frame}{\S3.4.3 Power-divergence family}



\begin{wideitemize}
\item
For multinomial data, both the Pearson statistic and deviance can be put under the umbrella of \emph{power-divergence statistics}  \cite{CressieRead1988}  given by
\[
S_\lambda =\sum_{i=1}^g \mbox{SD}_\lambda (\textbf{y}_i, \hat{\mbox{\boldmath $\pi$}}_i)
\]
where  
\begin{itemize}
\item
$
\mbox{SD}_\lambda (\textbf{y}_i, \hat{\mbox{\boldmath $\pi$}}_i)=\frac{2n_i}{\lambda (\lambda+1)}
\sum_{j=1}^k y_{ij}\left[\left(\frac{y_{ij}}{\hat{\pi}_{ij}}\right)^\lambda -1 \right]
$
\item 
$\lambda\in \mathbb{R}$ is a user chosen parameter.
\end{itemize}
\item The form of each $\mbox{SD}_\lambda (\textbf{y}_i, \hat{\mbox{\boldmath $\pi$}}_i)$ assesses the fit of the model by comparing the \emph{expected} and \emph{observed} frequencies.
\end{wideitemize}
\end{frame}


\begin{frame}{\S3.4.3 Power-divergence family (2)}
\begin{enumerate}
\item
At $\lambda=1$, $\displaystyle S_1=\sum_{i=1}^g n_i\sum_{j=1}^k \frac{(y_{ij}-\hat{\pi}_{ij})^2}{\hat{\pi}_{ij}}$,
the usual Pearson statistic.
\item
As $\lambda\rightarrow 0$, $\displaystyle S_0=2\sum_{i=1}^g n_i\sum_{j=1}^k y_{ij}
\log \frac{y_{ij}}{\hat{\pi}_{ij}}$,
the likelihood ratio statistic.
\item
At $\lambda\rightarrow -1$, $\displaystyle S_{-1}=2\sum_{i=1}^g n_i\sum_{j=1}^k \hat{\pi}_{ij}
\log \frac{\hat{\pi}_{ij}}{y_{ij}}$, Kullback's minimum discrimination information statistic.
\item
At $\lambda= -2$, 
$\displaystyle S_{-2}=\sum_{i=1}^g n_i\sum_{j=1}^k \frac{(y_{ij}-\hat{\pi}_{ij})^2}{y_{ij}}$,
Neyman's minimum modified $\chi^2$-statistic.
\item
At $\lambda=-\frac{1}{2}$, $\displaystyle S_{-\frac{1}{2}}\!=\!4\!\sum_{i=1}^g \sum_{j=1}^k 
(\!\sqrt{n_iy_{ij}}-\sqrt{n_i\hat{\pi}_{ij}})^2$, {\small Freeman-Tukey statistic.}
\end{enumerate}
In general, $\lambda\in [-1, 2]$ is recommended.
\end{frame}

\begin{frame}{\large \S3.4.3 Classical ``fixed cells" asymptotics}
Classical asymptotics for grouped data:
\begin{enumerate}
\item
$g$  fixed
\item
increasing  $n_i\rightarrow\infty$, with $\frac{n_i}{n}\rightarrow\lambda_i$
for fixed proportions $\lambda_i>0$, $i=1,\cdots,g$.
\item
a fixed ``number of cells" $k$ in each group
\item
a fixed  $p$ number of parameters.
\end{enumerate}
If these hold and the null is correct, then
$$S_\lambda \stackrel{a}{\sim} \chi^2 (g(k-1)-p). $$
Under local alternative hypothesis,  $S_\lambda$ converges to non-central $\chi^2$.
\end{frame}

\begin{frame}{\large \S3.4.3 sparseness and ``increasing-cells" asypmtotics}
\begin{wideitemize}
\item
The  number of observations
for a fixed covariate value $\textbf{x}$ is often small and the usual asymptotics will fail.
\item
Under such sparseness conditions, one can consider the asymptotic regime where both $n$ and $g$ grow to infinity.
(the ``increasing-cells" setting)
\item
$S_\lambda$ is no longer asymptotically $\chi^2$, but is  asymptotically
normal  under further regularity conditions. 

\item \emph{Personal note}: I have witnessed this phenomenon (``Chi-square becomes normal in high dimensions") many times,   in the currently active research area of {\bf high-dimensional statistics}. 
\end{wideitemize}
\end{frame}


\begin{frame}[allowframebreaks]
  \frametitle{References}    
  \begin{thebibliography}{10}    

%   \beamertemplatearticlebibitems
%  \bibitem[Lang, 1990]{Lang2014}
%    Lang, Joseph B
%    \newblock The Pearson Score Statistic for Multinomial- Poisson Models
%    \newblock \emph{Communications in Statistics-Theory and Methods 43.21 (2014): 4471-4491}
%  \beamertemplatearticlebibitems
%  \bibitem[Smyth, 2003 ]{smyth2003pearson}
%    Smyth, Gordon K
%    \newblock Pearson's goodness of fit statistic as a score test statistic
%    \newblock {\em Lecture notes-monograph series (2003): 115-126.}
%    
%      \beamertemplatearticlebibitems
%  \bibitem[Moore, 1977 ]{moore}
%    Moore David S
%    \newblock Generalized inverses, Wald's method, and the construction of chi-squared tests of fit
%    \newblock {\emph Journal of the American Statistical Association 72.357 (1977): 131-137}
%    
    
            \beamertemplatearticlebibitems
  \bibitem[Seber and Lee, 2003 ]{seberleelinear}
   Seber, G. A., \& Lee, A. J.
    \newblock Linear regression analysis, Second edition.
    \newblock { John Wiley \& Sons. 2003. }
    
    
            \beamertemplatearticlebibitems
  \bibitem[Rao and Mitra, 1972 ]{radhakrishna1972generalized}
 Radhakrishna Rao, C., and Sujit Kumar Mitra
    \newblock Generalized inverse of a matrix and its applications
    \newblock {In Proceedings of the Berkeley Symposium on Mathematical Statistics and Probability, vol. 1, pp. 601-620. University of California Press, Berkeley, 1972. }
    
       \beamertemplatearticlebibitems
      \bibitem[Tutz, 2012]{tutz2012}
   Tutz, Gerhard
    \newblock Regression for Categorical Data.
    \newblock {Cambridge 2012}
    
    
                \beamertemplatearticlebibitems
  \bibitem[Cressie and Read, 1988]{CressieRead1988}
 Noel A.C. Cressie, and Timothy R.C. Read
    \newblock Goodness-of-Fit Statistics for Discrete Multivariate Data
    \newblock {Springer-Verlag, 1988. }
  \end{thebibliography}
\end{frame}

\end{document}

